\documentclass[12pt,a4paper]{article}

\usepackage[T2A]{fontenc}
\usepackage[utf8]{inputenc}
\usepackage[russian]{babel}
\usepackage{geometry}
\usepackage{booktabs}
\usepackage{longtable}
\usepackage{array}
\usepackage{enumitem}
\usepackage{hyperref}
\usepackage{titlesec}
\usepackage{float}
\usepackage{amsmath}
\usepackage{amssymb}
\usepackage{multicol}
\usepackage{xcolor}

\geometry{left=2cm, right=1.5cm, top=2cm, bottom=2cm}

\hypersetup{
  colorlinks=true,
  linkcolor=black,
  urlcolor=blue,
  pdftitle={Методика преподавания математики и машинного обучения для 1-9 классов},
  pdfauthor={Аналитический отдел}
}

\titleformat{\section}{\Large\bfseries}{\thesection}{1em}{}
\titleformat{\subsection}{\large\bfseries}{\thesubsection}{1em}{}
\titleformat{\subsubsection}{\normalsize\bfseries}{\thesubsubsection}{1em}{}

\setlength{\parindent}{1cm}
\setlength{\parskip}{0.4em}
\setlength{\columnsep}{0.6cm}

\setlist[itemize]{noitemsep, topsep=2pt}
\setlist[enumerate]{noitemsep, topsep=2pt}

\newcolumntype{L}[1]{>{\raggedright\arraybackslash}p{#1}}
\newcolumntype{C}[1]{>{\centering\arraybackslash}p{#1}}
\newcolumntype{R}[1]{>{\raggedleft\arraybackslash}p{#1}}

\newcommand{\topic}[1]{\textbf{#1}}
\newcommand{\ml}[1]{\textcolor{blue!65!black}{\small #1}}

\newcommand{\note}[1]{%
  \begin{center}
  \fbox{\parbox{0.92\textwidth}{\small\textbf{Методическая заметка.} #1}}
  \end{center}
}

\title{
  \vspace{-1.2cm}
  \textbf{\Large МЕТОДИКА ПРЕПОДАВАНИЯ}\\[0.3cm]
  \textbf{\Large МАТЕМАТИКИ И МАШИННОГО ОБУЧЕНИЯ}\\[0.3cm]
  \textbf{\Large ДЛЯ УЧАЩИХСЯ 1--9 КЛАССОВ}\\[0.5cm]
  \rule{0.75\textwidth}{0.8pt}\\[0.4cm]
  \large Поурочный план на три года\\[0.2cm]
  \large (три занятия в неделю)\\[0.4cm]
  \rule{0.75\textwidth}{0.8pt}\\[0.5cm]
  \normalsize
  \begin{tabular}{rl}
    Младшие классы (1--4): & 1 час / занятие\\
    Средние классы (5--7): & 1{,}5 часа / занятие\\
    Старшие классы (8--9): & 1{,}5 часа / занятие\\
  \end{tabular}
}

\author{Аналитический отдел}
\date{Сентябрь 2026 г.}

\begin{document}

\maketitle

\vspace{1.2cm}

\begin{center}
\begin{minipage}{0.85\textwidth}
\small
\textbf{Аннотация.} Документ представляет собой полную методику преподавания
математики и машинного обучения (ML) учащимся 1--9 классов на горизонте трёх лет.
Содержит: педагогические принципы, сквозную логику «математика $\rightarrow$ ML»,
раздельные траектории для трёх возрастных групп, поурочные планы
всех 972 занятий (3 года $\times$ 3 траектории $\times$ 36 недель
$\times$ 3 занятия; по 324 занятия на траекторию), систему оценивания,
методику работы с одарёнными и отстающими учениками, а также чек-листы преподавателя.

\textbf{Ключевая идея.} Математика и ML преподаются не как два отдельных предмета,
а как единая система: математика 5--9 классов (комбинаторика, теория вероятностей,
статистика, алгебра) является прямым математическим фундаментом машинного обучения.
Такая связка даёт конкурентное преимущество: конкуренты либо учат только математику,
либо только программирование, но почти никто не строит явный мост между ними.
\end{minipage}
\end{center}

\vfill

\begin{center}
\small
\textbf{Структура документа}\\[0.3cm]
\begin{tabular}{rl}
1. & Педагогические принципы и нормативы\\
2. & Сквозная логика «математика $\rightarrow$ ML»\\
3. & Методика преподавания математики\\
4. & Три возрастные траектории: обзор\\
5. & Траектория А. Младшие классы (1--4), 1 час\\
6. & Траектория Б. Средние классы (5--7), 1{,}5 часа\\
7. & Траектория В. Старшие классы (8--9), 1{,}5 часа\\
8. & Система оценивания и аттестации\\
9. & Работа с одарёнными и отстающими\\
10. & Методика, снижающая отток учеников\\
11. & Типовые ошибки и как их предотвращать\\
12. & Чек-листы преподавателя\\
13. & Приложения: шаблоны, конспекты, материалы\\
\end{tabular}
\end{center}

\newpage
\tableofcontents
\newpage

% ====================================================================
\section{Педагогические принципы и нормативы}

\subsection{Десять принципов программы}

\begin{enumerate}
  \item \textbf{Принцип «сначала понимание, потом код».} Ученик не пишет программу,
  пока не может объяснить на словах, что она должна делать. Это снимает страх перед
  синтаксисом и переносит фокус на логику, которая и является математикой.
  \item \textbf{Математика --- это язык данных.} Каждая ключевая тема математики
  показывается в двух ипостасях: как инструмент решения задачи и как способ описать
  данные для модели.
  \item \textbf{ML как «обучение на примерах», а не как магия.} На каждом этапе
  термин «обучение» раскрывается через понятную аналогию (детектор спама,
  рекомендатель книг), прежде чем появляется настоящий код.
  \item \textbf{Видимый результат в конце каждого занятия.} Ученик должен уйти
  с ощущением «я сегодня что-то построил/доказал/получил», а не «мы послушали теорию».
  \item \textbf{Ошибка --- часть обучения.} Каждое занятие включает
  обязательный этап разбора ошибок, без осуждения.
  \item \textbf{Спиральное повторение.} Ключевые темы возвращаются в следующих
  классах на новом уровне (дроби $\rightarrow$ вероятность $\rightarrow$ метрики).
  \item \textbf{Минимальная теория перед максимумом практики.} Не более
  15--20 минут чистого объяснения за занятие.
  \item \textbf{Олимпиадный элемент с 3 класса.} Даже для младших учеников
  предусмотрены задачи уровня начальной олимпиады --- это создаёт
  «потолок роста» и удерживает сильных учеников.
  \item \textbf{Проектная защита вместо экзамена.} Итогом каждого года
  является работа, которую ученик сам защищает.
  \item \textbf{Принцип «одна переменная за раз».} При отладке кода ученик
  проверяет по одной гипотезе за раз --- это переносит дисциплину эксперимента
  из программирования в математику и обратно.
\end{enumerate}

\subsection{Нормативы нагрузки}

\begin{table}[H]
\centering
\caption{Нормативная нагрузка по траекториям}
\label{tab:norms}
\begin{tabular}{L{3.2cm}C{2.6cm}C{2.6cm}C{2.4cm}C{2.4cm}}
\toprule
\textbf{Траектория} & \textbf{Занятий в неделю} & \textbf{Часов в неделю} &
\textbf{Занятий в год} & \textbf{Часов в год}\\
\midrule
А. Младшие (1--4) & 3 & 3,0 & 108 & 108\\
Б. Средние (5--7) & 3 & 4,5 & 108 & 162\\
В. Старшие (8--9) & 3 & 4,5 & 108 & 162\\
\bottomrule
\end{tabular}
\end{table}

\noindent
Расчёт: 36 учебных недель $\times$ 3 занятия = 108 занятий в год,
то есть по 324 занятия на траекторию за три года. В сумме по трём
траекториям --- 972 занятия. По времени: траектория А --- 324 часа
(1 час на занятие), траектории Б и В --- по 486 часов
(1{,}5 часа на занятие), всего 1296 часов.

\note{В расчёте 108 занятий в год заложено 36 учебных недель с трёхнедельным
каникулярным резервом внутри учебного года. Если в конкретном лицее
34 учебные недели, из каждого блока плана исключается одна резервная
тема (они помечены в плане знаком $\star$), а освободившееся занятие
используется для повторения или индивидуальной работы.}

\subsection{Целевые результаты по годам}

\begin{longtable}{L{2.4cm}L{4.3cm}L{4.3cm}L{4.3cm}}
\toprule
\textbf{Год} & \textbf{Математика} & \textbf{ML и программирование} & \textbf{Личностные}\\
\midrule
\endfirsthead

\toprule
\textbf{Год} & \textbf{Математика} & \textbf{ML и программирование} & \textbf{Личностные}\\
\midrule
\endhead

1 & Уверенный счёт в пределах таблицы умножения; понимание смысла дроби,
процента, отношения; решение логических и комбинаторных задач
начального уровня & Базовые навыки алгоритмического мышления: линейный
алгоритм, ветвление, цикл. Умение описать задачу словами и
закодировать её визуально. Понимание, что такое «данные» &
Умение доводить задачу до конца; адекватная реакция на ошибку;
готовность проверять чужое решение\\
\addlinespace
2 & Систематизация дробей, процентов, пропорций; введение в алгебру
(числовые и буквенные выражения, уравнения, неравенства); комбинаторика
основных правил; элементы геометрии (площади, объёмы, подобие) & Язык
Python: переменные, условия, циклы, функции, списки. Основы
визуализации данных. Построение первой модели и оценка её точности &
Работа в паре и в команде; навык декомпозиции большой задачи;
готовность признавать свои ошибки\\
\addlinespace
3 & Алгебра и планиметрия на уровне подготовки к ОГЭ/олимпиадам;
теория вероятностей и математическая статистика --- фундамент ML;
числовые и функциональные задачи & Реальные проекты на данных:
построение и обучение модели, анализ качества, интерпретация
результатов. Понимание переобучения и разделения выборки &
Способность к самостоятельному исследованию; умение защитить
решение; планирование собственного обучения\\
\bottomrule
\end{longtable}

% ====================================================================
\section{Сквозная логика «математика $\rightarrow$ ML»}

\subsection{Таблица соответствий}

Ключевая методическая опора всей программы --- явное указание того,
какая тема математики становится фундаментом для какой техники ML.
Эта таблица должна быть перед глазами преподавателя при планировании
каждого блока.

\begin{longtable}{L{3.4cm}L{4.6cm}L{4.6cm}L{2.6cm}}
\toprule
\textbf{Математика} & \textbf{Техника ML} & \textbf{Почему именно это} & \textbf{Классы}\\
\midrule
\endfirsthead

\toprule
\textbf{Математика} & \textbf{Техника ML} & \textbf{Почему именно это} & \textbf{Классы}\\
\midrule
\endhead

Линейные уравнения &
Регрессия (линейная модель) & Уравнение прямой $y = wx + b$ --- это
и есть линейная регрессия; ученик видит буквальное совпадение &
5--9\\

\addlinespace
Квадратные уравнения &
Полиномиальные признаки & Подстановка функции признака в модель
даёт нелинейность без изменения «механики» &
6--9\\

\addlinespace
Дроби, пропорции &
Нормировка и масштабирование & Отношение $x / \max(x)$ --- это
и есть нормировка; пропорция объясняет, зачем это нужно &
5--8\\

\addlinespace
Вектор и координаты &
Признаки (features) и их вектор & Набор признаков модели ---
это буквально вектор координат точки в пространстве признаков &
7--9\\

\addlinespace
Скалярное произведение &
Линейная комбинация признаков & $w \cdot x$ --- ровно то, что
вычисляет нейрон и линейная модель &
8--9\\

\addlinespace
Комбинаторика, правило
произведения &
Множество комбинаций признаков,
перебор гипотез & Число способов выбрать набор признаков растёт
как $2^n$ --- это объясняет «проклятие размерности» &
6--9\\

\addlinespace
Теория вероятностей &
Вероятность, распределения,
байесовская классификация & Прямой и практический мост: условная
вероятность $\rightarrow$ формула Байеса $\rightarrow$ спам-фильтр &
8--9\\

\addlinespace
Математическая
статистика &
Выборка, среднее, дисперсия,
доверительные интервалы & Оценка качества модели теми же
средствами, что и в школьной статистике &
7--9\\

\addlinespace
Логарифм &
Логарифмическая шкала, softmax &
Логарифм объясняет, почему ошибки «схлопываются», и подводит
к пониманию cross-entropy &
9\\

\addlinespace
Графики функций &
Визуализация и обучение кривой
(градиентный спуск) &
«Поднялись выше --- скатились ниже» --- это и есть градиентный
спуск на картинке &
8--9\\

\addlinespace
Модуль числа, расстояние &
Метрики качества (MAE, RMSE) &
Средняя абсолютная ошибка буквально является средним
модулем отклонения &
8--9\\

\addlinespace
Медиана &
Робастность модели, устойчивость
к выбросам &
Медиана не «прыгает» от одного выброса --- ровно за это
её ценят в ML &
8--9\\

\addlinespace
Множество, мощность &
Множество значений признака,
кодирование категорий & Терминология и логика совпадают; это
естественная точка входа в тему &
7--9\\
\bottomrule
\end{longtable}

\note{Эта таблица --- главный навигационный инструмент преподавателя.
Перед каждым новым блоком ML рекомендуется за 5 минут повторить
соответствующую тему математики. И наоборот: каждая новая тема
математики должна «выстреливать» в конкретную технику ML
в течение 2--3 занятий. Это создаёт ощущение единой системы,
а не двух разрозненных предметов.}

\subsection{Принцип «обучение на примерах», объяснённый через математику}

Ключевая педагогическая задача --- показать, что ML --- это не «магия
нейросетей», а статистика. Ниже --- готовая аналогия для каждого
возраста.

\begin{table}[H]
\centering
\caption{Аналогия обучения для разных возрастов}
\label{tab:analogy}
\begin{tabular}{L{1.6cm}L{6.4cm}L{7.2cm}}
\toprule
\textbf{Класс} & \textbf{Аналогия} & \textbf{Как используется на занятии}\\
\midrule
1--2 & «Учимся отличать кошку от собаки:
смотрим на картинки и запоминаем
что общего» & Учитель показывает карточки, просит ребёнка
объяснить правило. Затем --- закрепление на
новых карточках. Это интуитивное описание
обучения модели\\
\addlinespace
3--4 & «Различаем чётные и нечётные числа:
нашли закономерность --- правило» & Формулируется первое
«правило модели». Вводится термин
признак (чётность)\\
\addlinespace
5--6 & «Учитель ставит оценки за
контрольные по нашим ответам» &
Формулируется понятие обучающей выборки
(наши ответы) и проверки (оценка учителя)\\
\addlinespace
7 & «Строим график и смотрим, какая
прямая ближе всего к точкам» &
Прямой переход к линейной регрессии;
показывается переобучение: прямая,
проходящая через все точки, --- плохая\\
\addlinespace
8 & «Какая карточка счастливее?
Считаем среднее по выборке и сравниваем
средние» & Вводится среднее и дисперсия как
инструмент сравнения групп; именно они
лежат в основе дерева решений\\
\addlinespace
9 & «Кто сегодня чаще попадал в
аварию, сравним вероятности с учётом
погоды» & Формула Байеса на бытовых примерах,
затем --- спам-фильтр как реальная задача\\
\bottomrule
\end{tabular}
\end{table}

% ====================================================================
\newpage
\section{Методика преподавания математики}

Математика в этом курсе --- не «подготовка к ML» и не вспомогательный
предмет. Это полноценная дисциплина с собственными целями, логикой
и системой контроля. Настоящий раздел описывает, \emph{как именно}
она преподаётся: от целей и содержания до приёмов работы, типичных
ошибок учеников и форм проверки усвоения.

\subsection{Место математики в курсе}

\noindent
Математика занимает \textbf{60\,\%} учебного времени занятия и
\textbf{60\,\%} домашнего задания. Оставшиеся 40\,\% отданы
программированию и машинному обучению. Баланс поддерживается жёстко:
если на трёх занятиях подряд математика сведена к «разминке»,
программа считается нарушенной и тема возвращается в следующем блоке.

Математика выполняет три функции, и все три равноправны.

\begin{enumerate}
  \item \textbf{Инструментальная.} Математика даёт язык, на котором
  записывается машинное обучение. Без линейной алгебры, теории
  вероятностей и статистики техника ML остаётся набором рецептов,
  которые невозможно ни проверить, ни исправить.
  \item \textbf{Диагностическая.} Именно на математике проявляется
  непонимание ученика. Отсюда следует, что доля отстающих в курсе
  с сильным математическим блоком ниже: преподаватель видит проблему
  на школьном материале, а не на абстрактном коде.
  \item \textbf{Самостоятельная.} Значительная часть учеников готовится
  к инженерным и естественно-научным направлениям. Курс даёт им
  систематическую подготовку, соответствующую школьной программе,
  но с опорой на задачи повышенной сложности и раннее введение
  олимпиадных приёмов.
\end{enumerate}

\note{Математика в этом курсе никогда не подаётся как «средство для
ML». Ученик должен понимать, что математика ценна сама по себе, иначе
связь перестаёт работать: связь возникает между двумя самостоятельными
предметами, а не внутри одного.}

\subsection{Десять принципов преподавания математики}

\begin{enumerate}
  \item \textbf{Сначала смысл, потом символ.} Ни одна формула не
  вводится раньше, чем ученик увидел задачу, которую она решает.
  Формула --- это сжатая запись уже понятого действия.
  \item \textbf{Три ступени представления.} Каждая тема проходит
  путь «предметная $\rightarrow$ наглядная $\rightarrow$ символическая».
  Переход на следующую ступень без освоения предыдущей запрещён.
  \item \textbf{Одна тема --- одна задача-ситуация.} Новое понятие
  всегда рождается из конкретной задачи, которую невозможно решить
  известными способами.
  \item \textbf{Показывать границы применимости.} Вместе с правилом
  обязательно проговаривается, когда оно не работает. Ученик,
  знающий границы, не применяет инструмент неверно.
  \item \textbf{Устная работа обязательна.} Не менее 5 минут устного
  счёта или устного рассуждения каждое занятие --- это тренирует
  переход от черновика к голове.
  \item \textbf{Рассуждение важнее ответа.} Задача без записанного
  рассуждения считается нерешённой, даже если ответ верен.
  \item \textbf{Ошибку разбираем, а не отмечаем.} Разбор ошибок
  входит в хронометраж каждого занятия и не связан с выставлением
  оценок.
  \item \textbf{Двойная запись.} Каждая тема записывается дважды:
  как решение задачи и как описание данных для модели. Это делает
  таблицу соответствий из раздела~2 рабочим инструментом, а не
  декларацией.
  \item \textbf{Спираль с разной глубиной.} Ключевая тема возвращается
  в каждом году --- дроби как понятие, дроби как вероятность, дроби
  как нормировка, --- но на качественно новом уровне.
  \item \textbf{Проверка усвоения, а не посещаемости.} Процент
  выполненных домашних заданий не является показателем освоения.
  Показатель --- доля задач, решённых самостоятельно на оценочном занятии.
\end{enumerate}

\subsection{Содержание математики по классам}

Ниже --- целевое содержание математического блока. Таблица нужна
преподавателю как основа сверки: каждая тема школьной программы
должна быть закрыта, а каждая закрытая тема --- «выстрелить» в
конкретную технику ML в течение двух-трёх занятий.

\begin{longtable}{L{1.1cm}L{4.4cm}L{4.4cm}L{4.4cm}}
\toprule
\textbf{Кл.} & \textbf{Разделы математики} & \textbf{Ключевые понятия
и навыки} & \textbf{Мост в ML}\\
\midrule
\endfirsthead

\toprule
\textbf{Кл.} & \textbf{Разделы математики} & \textbf{Ключевые понятия
и навыки} & \textbf{Мост в ML}\\
\midrule
\endhead

1 & Числа до 100, сложение и вычитание, простые фигуры, длина,
время, задачи в 1--2 действия & Счёт предметов, сравнение,
«больше --- меньше», геометрические фигуры & Отсчёт объектов;
признак «как посчитать»\\
\addlinespace
2 & Числа до 1000, таблица умножения, умножение и деление, углы,
периметр, площадь, длина--масса--вместимость & Кратные числа,
разложение на множители, единицы измерения & Признак
«кратность признака», нормализация единиц\\
\addlinespace
3 & Умножение в столбик, деление с остатком, доли, координаты
на прямой, окружность и круг, площадь фигур & Доля, часть и целое,
координата, длина окружности & Признак как доля признака;
окружность как граница множества\\
\addlinespace
4 & Многозначные числа, деление углом, дроби и смешанные числа,
площади, движение, координаты & Понятие дроби, полный квадрат,
время--скорость--расстояние & Признак как дробная величина;
нормировка $x / \max(x)$\\
\midrule
5 & Десятичные дроби и действия с ними, проценты, дроби,
начала алгебры, треугольник и прямоугольник, объёмы, координаты
на плоскости & Процент как дробь, перевод дроби в десятичную,
формула, объём, отрицательные координаты & Нормировка признаков;
линейная зависимость; формула как функция\\
\addlinespace
6 & Рациональные числа, отношения и пропорции, масштаб, решение
задач на движение, окружность и круг, симметрия & Пропорция, коэффициент
подобия, скорость, положительные и отрицательные числа & Пропорция как
масштабирование; обратная пропорция; знак признака\\
\midrule
7 & Линейные уравнения, одночлены и многочлены, формулы
сокращённого умножения, график и линейная функция, степени,
треугольник, окружность & Уравнение как равенство, коэффициенты,
угловой коэффициент, свойство функции & \textbf{Линейная регрессия}
прямо из уравнения $y = wx + b$\\
\addlinespace
8 & Квадратные уравнения, иррациональные числа, корни, неравенства,
функции и графики, площади, подобные треугольники & Дискриминант,
разделение множеств, подобие, площадь & \textbf{Полиномиальные признаки};
решение неравенств; метрики ошибки\\
\addlinespace
9 & Рациональные выражения, прогрессии, системы уравнений, векторы
и координаты, скалярное произведение, вероятность и статистика &
Вектор, скалярное произведение, комбинаторика, условная вероятность,
выборка, среднее, дисперсия & \textbf{Веса признаков}, вектор,
скалярное произведение $w \cdot x$, \textbf{формула Байеса},
спам-фильтр, среднее и дисперсия как метрики\\
\bottomrule
\end{longtable}

\note{Столбец «Мост в ML» --- не цель, а проверка методики. Если тему
не удаётся «выстрелить», это не значит, что она лишняя: значит, что
следует искать связь в более общем виде (масштабирование, сравнение
групп, подсчёт вариантов). Темы без естественной связи берутся
как самостоятельный фундамент и в таком случае учатся
с обязательной олимпиадной составляющей.}

\subsection{Методика введения новой темы}

Введение любой темы состоит из пяти обязательных этапов. Этапы
не переставляются и не сокращаются: если ученик не освоил этап~2,
этап 3 не запускается.

\begin{table}[H]
\centering
\caption{Пять этапов введения темы}
\label{tab:newtopic}
\begin{tabular}{L{1.4cm}L{3.4cm}L{5.0cm}L{4.2cm}}
\toprule
\textbf{Этап} & \textbf{Название} & \textbf{Что происходит} &
\textbf{Признак готовности}\\
\midrule
1 & Ситуация & Ученик получает задачу, которую не может решить
известным способом, и сам формулирует, чего ему не хватает &
Ученик словами назвал недостающее знание\\
\addlinespace
2 & Предметная модель & Проблема решается с реальными предметами
или наглядными схемами: полоска бумаги, фишки, монеты, линейка &
Решение получено, схема нарисована\\
\addlinespace
3 & Символическая запись & То же самое записывается знаками:
формулой, выражением, схемой решения & Запись воспроизводится
самостоятельно\\
\addlinespace
4 & Примеры и контрпримеры & Тема проверяется на 3--4 примерах,
подбирается пример, где приём не работает &
Ученик сам находит границу применимости\\
\addlinespace
5 & Применение & Тема применяется в задачах всех трёх уровней
сложности и связывается с техникой ML &
Есть собственная задача и связь с ML\\
\bottomrule
\end{tabular}
\end{table}

\noindent
Этап 1 занимает 3--5 минут и проводится в устной форме. Этапы 2 и 3
занимают основную часть занятия. Этапы 4 и 5 частично выносятся
в домашнее задание.

\subsection{Устный счёт и математическая грамотность}

Устная работа --- обязательный элемент, а не разминка «для галочки».
Она решает три задачи: тренирует автоматизм, вскрывает ошибки
в устной речи и готовит к олимпиадным ситуациям, где считать
некогда.

\begin{table}[H]
\centering
\caption{Форматы устной работы}
\label{tab:oral}
\begin{tabular}{L{3.2cm}C{1.8cm}L{5.0cm}L{4.0cm}}
\toprule
\textbf{Формат} & \textbf{Мин} & \textbf{Содержание} &
\textbf{Тренирует}\\
\midrule
Цепочка & 3 & Ученик продолжает цепочку: $3 \times 4$, $\times 5$,
$\times 6$ и так далее & Умножение в уме, удержание
результата\\
\addlinespace
Обратный счёт & 2 & Вычитание или деление с уменьшением
последовательности & Устойчивость к ошибке\\
\addlinespace
Случайное условие & 3 & Учитель меняет числа: «а теперь
другой множитель» & Гибкость, отказ от заученного шаблона\\
\addlinespace
Прикидка & 3 & «Округлите и оцените порядок ответа» &
Оценка правдоподобности, защита от грубых ошибок\\
\addlinespace
Быстрая модель & 3 & Оценка задачи без точного решения &
Моделирование, связь с ML-предсказанием\\
\bottomrule
\end{tabular}
\end{table}

\note{Приём «прикидка» --- прямой мост к машинному обучению.
Прежде чем обучать модель, ученик учится сам оценивать правдоподобность
ответа. Умение сказать «ответ должен быть примерно 300, а я получил
3000 --- значит, ошибка в расчёте» работает и в арифметике,
и в проверке вывода модели.}

\subsection{Работа с доказательствами}

Цель --- не научить доказывать по образцу, а сформировать привычку
рассуждать и записывать рассуждение.

\begin{enumerate}
  \item \textbf{Учим на коротких утверждениях.} Первые доказательства
  строятся на 3--4 предложениях: «чётное число имеет вид $2n$»,
  «сумма двух чётных чётна».
  \item \textbf{Шаблон рассуждения обязателен.} Ученик проговаривает
  структуру: «Рассмотрим произвольный... Покажем, что... Значит...».
  Пустая форма, заполняемая учеником, используется как
  отправная точка.
  \item \textbf{Разбор образцов, а не диктовка.} Преподаватель
  разбирает 2--3 готовых доказательства, затем ученики
  самостоятельно находят ошибку в неверном доказательстве.
  \item \textbf{Градация сложности.} Олимпиадные доказательства
  даются в младших классах только в устной форме; запись
  обязательна с 5 класса.
\end{enumerate}

\subsection{Типичные ошибки учеников и их профилактика}

Здесь перечислены ошибки \emph{учеников} в отличие от раздела~11,
где разбираются ошибки преподавателя.

\begin{longtable}{L{4.6cm}L{4.6cm}L{5.0cm}}
\toprule
\textbf{Ошибка} & \textbf{Почему возникает} & \textbf{Профилактика}\\
\midrule
\endfirsthead

\toprule
\textbf{Ошибка} & \textbf{Почему возникает} & \textbf{Профилактика}\\
\midrule
\endhead

Потеря единиц измерения & Нет устойчивой привычки
подписывать единицы & Запись единиц --- обязательное условие
сдачи задачи; 1 минута на единицы в каждом ДЗ\\
\addlinespace
Перенос свойства с одного объекта на весь класс & Формальное
заучивание правила без контрпримера & Обязательный этап 4:
подобрать пример, где правило не работает\\
\addlinespace
Действие по образцу, хотя условие изменилось & Ученик
узнал тип задачи, но не прочитал условие & Приём «прочитай
условие вслух и назови, что просят»\\
\addlinespace
Путаница «больше --- меньше» в неравенствах & Направление
умножения меняется вместе со знаком & Визуальная проверка:
отметить точку на числовой прямой до записи ответа\\
\addlinespace
Считают, что $x^2 = 2x$ & Степень и умножение сливаются
в одну операцию & Обязательный устный вопрос: «что здесь
умножается и сколько раз?»\\
\addlinespace
Ошибка в знаке при раскрытии скобок & Слепая запись по шаблону
без контроля & Приём «последняя операция первой»:
раскрывать справа налево\\
\addlinespace
Делимость проверяется «на глаз» & Нет алгоритма проверки &
Вводится явный алгоритм проверки признака делимости\\
\addlinespace
Вероятность больше 1 или отрицательная & Вероятность
воспринимается как «доля ответа», а не множества &
Обязательная числовая проверка $0 \leq p \leq 1$
с объяснением, что она означает\\
\addlinespace
Решение найдено, но не проверено & Экономия времени
на этапе проверки & Правило «обязательной проверки»:
подстановка в условие — 2 минуты\\
\addlinespace
Прикидка не делается & Ответ выглядит правдоподобно
независимо от вычислений & Каждый ответ сверяется
с прикидкой, сделанной до вычисления\\
\bottomrule
\end{longtable}

\subsection{Олимпиадная составляющая}

Математика сохраняет олимпиадную линию во всех траекториях:
это единственный способ создать «потолок роста» для сильного
ученика, не увеличивая нагрузку на остальных.

\begin{table}[H]
\centering
\caption{Олимпиадная линия по траекториям}
\label{tab:olymp}
\begin{tabular}{L{3.0cm}C{2.2cm}L{4.8cm}L{4.0cm}}
\toprule
\textbf{Траектория} & \textbf{Формат} & \textbf{Приёмы} &
\textbf{Где в плане}\\
\midrule
А. Младшие & Задачи-вызовы без формул, устное рассуждение,
мини-олимпиада раз в 4 недели & Принцип Дирихле, раскраска,
комбинаторика, логика & Занятия с $\star$, блок 3\\
\addlinespace
Б. Средние & Олимпиадные задачи уровня школы, тренировочные
туры по 3 задачи & Принцип крайнего, инвариант, разбор случаев,
принцип Дирихле & Занятия с $\star$, блоки 2 и 3\\
\addlinespace
В. Старшие & Муниципальный этап ВсОШ, олимпиадные туры
по 5--6 задач & Графы, инвариант, математическая индукция,
неравенства & Блоки 2 и 3, занятия 73--108\\
\bottomrule
\end{tabular}
\end{table}

\note{Правило допуска к олимпиадной задаче: ученик имеет право
не знать решение, но обязан записать ход рассуждения --- даже
если он пришёл к неверному ответу. Записано рассуждение --- работа
засчитана, это главный навык олимпиады.}

\subsection{Проверка усвоения математики}

Оценивание математики строится на четырёх составляющих. Ни одна
из них не сводится к количеству сделанных заданий.

\begin{enumerate}
  \item \textbf{Теоретическое знание} (30\,\%): умение объяснить
  понятие своими словами, назвать единицы измерения, привести
  пример и контрпример, указать границы применимости.
  \item \textbf{Техника вычислений} (30\,\%): точность и скорость
  счёта, выполнение алгоритмов без ошибок подстановки.
  \item \textbf{Рассуждение} (25\,\%): корректный ход решения
  в задачах с неполным условием, обоснование выбора метода.
  \item \textbf{Практика и применение} (15\,\%): умение решать
  бытовую, финансовую и геометрическую задачу целиком,
  включая прикидку и проверку ответа.
\end{enumerate}

\noindent
Оценочные мероприятия по математике: устный счёт на каждом занятии,
три уровня домашних заданий, мини-олимпиада раз в четыре недели,
итоговая работа в конце каждого блока из 36 занятий и защита
математической части проекта.

\subsection{Математические проекты}

Проект --- обязательный элемент каждого года по математике
независимо от ML-проекта. Тема выбирается учеником из списка
и защищается устно.

\begin{table}[H]
\centering
\caption{Математические проекты по траекториям}
\label{tab:mathproj}
\begin{tabular}{L{3.0cm}L{4.6cm}L{4.4cm}L{2.4cm}}
\toprule
\textbf{Траектория} & \textbf{Проект} & \textbf{Математическое
содержание} & \textbf{Защита}\\
\midrule
А. Младшие & «Модель моего класса» & Сбор данных о росте,
построение диаграммы, среднее и сравнение групп & 3--5 мин\\
\addlinespace
Б. Средние & «Бюджет класса» & Проценты, пропорции, округление,
среднее и медиана расходов & 5--7 мин\\
\addlinespace
Б. Средние & «Календарь и вероятность» & Комбинаторика,
правило произведения, оценка вероятности & 5--7 мин\\
\addlinespace
В. Старшие & «Модель успеваемости» & Линейная регрессия,
коэффициент корреляции, прогноз & 7--10 мин\\
\addlinespace
В. Старшие & «Справедливость правила» & Байесова классификация,
ложные срабатывания, критерий Байеса & 7--10 мин\\
\bottomrule
\end{tabular}
\end{table}

\note{Защита проекта по математике оценивается по трём критериям:
корректность математической части (60\,\%), полнота рассуждения
(25\,\%), ответы на вопросы (15\,\%). Оценка за красивую презентацию
не выставляется.}

% ====================================================================
\section{Три возрастные траектории: обзор}

\subsection{Сравнительная карта программ}

\begin{longtable}{L{3.2cm}L{3.7cm}L{3.7cm}L{3.7cm}}
\toprule
\textbf{Параметр} & \textbf{А. Младшие (1--4)} &
\textbf{Б. Средние (5--7)} & \textbf{В. Старшие (8--9)}\\
\midrule
\endfirsthead

\toprule
\textbf{Параметр} & \textbf{А. Младшие (1--4)} &
\textbf{Б. Средние (5--7)} & \textbf{В. Старшие (8--9)}\\
\midrule
\endhead

Длительность & 60 мин & 90 мин & 90 мин\\
\addlinespace
Ядро технологии &
Визуальное программирование
(Scratch, блочные среды) &
Текстовый Python в браузере &
Python + pandas + scikit-learn\\
\addlinespace
Математика & Арифметика, логика,
комбинаторика &
Дроби, проценты, начала алгебры &
Алгебра, вероятность, статистика\\
\addlinespace
Уровень ML & «Модель как правило»,
сортировка & Первая обученная модель,
линейная регрессия & Классификация, регрессия,
нейросети, переобучение\\
\addlinespace
Итог года & Игра-проект в Scratch &
Анализ реального набора данных &
Собственный ML-проект с защитой\\
\addlinespace
Главный риск &
Играть, а не учиться &
Упирание в синтаксис &
Страх «сложной математики»\\
\bottomrule
\end{longtable}

\subsection{Структура каждого учебного года}

Применяется единая сетка --- 36 учебных недель, разделённых на
три блока по 12 недель (по 36 занятий в каждом).

\begin{table}[H]
\centering
\caption{Сетка учебного года}
\label{tab:yeargrid}
\begin{tabular}{L{3.4cm}C{2.6cm}L{6.2cm}}
\toprule
\textbf{Блок} & \textbf{Занятия} & \textbf{Содержание}\\
\midrule
Блок 1 & 1--36 & Новый материал + закрепление\\
Блок 2 & 37--72 & Углубление + первые проекты\\
Блок 3 & 73--108 & Практика, олимпиады, итоговый проект\\
\bottomrule
\end{tabular}
\end{table}

\subsection{Шаблоны занятий}

\subsubsection*{Занятие 60 минут (траектория А)}

\begin{table}[H]
\centering
\caption{Хронометраж 60-минутного занятия}
\begin{tabular}{L{2.2cm}C{1.8cm}L{9.2cm}}
\toprule
\textbf{Этап} & \textbf{Мин} & \textbf{Содержание}\\
\midrule
Вход & 5 & Разминка: устный счёт / повторение вчерашнего\\
Теория & 10 & Один новый concept. Объяснение через аналогию\\
Демонстрация & 7 & Учитель показывает на экране или на бумаге\\
Практика & 25 & Ученики решают сами, в парах. Учитель --- помощник\\
Мини-итог & 8 & Задача-вызов или разбор ошибок\\
Финал & 5 & Что узнали, что дальше, домашнее задание\\
\bottomrule
\end{tabular}
\end{table}

\subsubsection*{Занятие 90 минут (траектории Б и В)}

\begin{table}[H]
\centering
\caption{Хронометраж 90-минутного занятия}
\begin{tabular}{L{2.4cm}C{1.6cm}L{9.2cm}}
\toprule
\textbf{Этап} & \textbf{Мин} & \textbf{Содержание}\\
\midrule
Вход & 5 & Разминка, проверка домашнего задания\\
Часть I. Теория & 12 & Первый новый concept\\
Часть I. Практика & 18 & Упражнения с нарастающей сложностью\\
Перерыв & 5 &\\
Часть II. Теория & 12 & Второй concept, часто --- мост к ML\\
Часть II. Практика & 20 & Работа в парах, мини-проект\\
Разбор & 8 & Ошибки, разбор олимпиадных задач\\
Финал & 5 & Итоги, домашнее задание, анонс\\
\addlinespace
\multicolumn{2}{r}{\textbf{Резерв}} & \textbf{5 мин --- на «добор» или обсуждение}\\
\bottomrule
\end{tabular}
\end{table}

\note{Правило хронометража жёсткое: теория никогда не занимает
больше 24 минут за занятие. Если тема не уложилась --- переносите
на следующее занятие, но не сокращайте практику. Ученик
запоминает не теорию, а тот опыт, который получил своими руками.}

\subsection{Система домашних заданий}

\begin{table}[H]
\centering
\caption{Нормы домашних заданий}
\label{tab:hw}
\begin{tabular}{L{3.6cm}L{2.4cm}L{3.4cm}L{4.2cm}}
\toprule
\textbf{Траектория} & \textbf{Мин на ДЗ} & \textbf{Объём} & \textbf{Форма сдачи}\\
\midrule
А. Младшие & 15--20 & 2--3 задачи & Устный ответ на следующем
занятии или тетрадь\\
Б. Средние & 25--35 & 4--5 задач + код & Тетрадь + файл\\
В. Старшие & 40--50 & 6--8 задач + код & Тетрадь + код + README\\
\bottomrule
\end{tabular}
\end{table}

\noindent
\textbf{Правило 60/40:} 60\% домашнего задания --- математика,
40\% --- программирование и ML. Это сохраняет баланс и не даёт
технике вытеснить фундамент.

% ====================================================================
\newpage
\section{Траектория А. Младшие классы (1--4), 1 час}

\subsection{Общие положения траектории}

\begin{itemize}
  \item \textbf{Возраст:} 7--11 лет. Длительность занятия --- 60 минут.
  \item \textbf{Главная задача года:} не программирование, а
  \emph{алгоритмическое мышление} через блочную среду и логические задачи.
  \item \textbf{Среда:} Scratch (оффлайн-версия для офлайн-заний, онлайн --- для дома).
  \item \textbf{Особенность:} у младших учеников нет школьного интереса к
  «настоящей математике», поэтому математика подаётся через игры,
  соревнования и наглядные модели. Текстовый код --- только с 5 класса.
  \item \textbf{Методический риск:} увлечение созданием игр без понимания
  логики. Решается правилом: «каждая игра должна быть объяснена словами
  перед сборкой».
\end{itemize}

\subsection{Год 1. Тематический план}

\begin{longtable}{C{1.0cm}L{4.4cm}L{4.0cm}L{4.4cm}}
\toprule
\textbf{№} & \textbf{Тема} & \textbf{Математика} & \textbf{Практика (Scratch/ML)}\\
\midrule
\endfirsthead

\toprule
\textbf{№} & \textbf{Тема} & \textbf{Математика} & \textbf{Практика (Scratch/ML)}\\
\midrule
\endhead

\multicolumn{4}{l}{\textbf{Блок 1. Мыслим как программисты (занятия 1--36)}}\\
\midrule
1 & Что такое программа & Порядок действий, следование инструкции & Инструкция для настольной игры\\
2 & Линейный алгоритм & Числовые действия по шагам & Блоки «повтори», «сказать»\\
3 & Исполнитель & Кто и что выполняет & Спрайт как исполнитель\\
4 & Блок-схема & Представление схемы & Собираем первую блок-схему\\
5 & Разветвление & Сравнение чисел, знаки & Условие «если---то---иначе»\\
6 & Истина и ложь & Логические высказывания & Проверка условий в Scratch\\
7 & Цикл & Повторение действий & Цикл «повторить N раз»\\
8 & Вложенный цикл & Счёт в нескольких кругах & Рисуем сетку кругами\\
9 $\star$ & Итог блока 1 & Повторение & Мини-игра за 5 минут\\
10 & Переменные & Что такое «величина» & Создание переменной, счётчик очков\\
11 & Списки & Хранение множества & Массив значений в Scratch\\
12 & Условие + список & Выбор из коллекции & «Кто сегодня дежурит»\\
13 & Ввод данных & Как задать вопрос и получить ответ & Блок «спросить и ждать»\\
14 & Математика в цикле & Сумма чисел от 1 до N & Вычисляем сумму циклом\\
15 & Разряды & Разложение числа на разряды & Считаем цифры числа\\
16 $\star$ & Мини-проект 1 & --- & Игра «Угадай число»\\
17 & Больше, меньше, равно & Сравнение, знаки & Ветвление в игре\\
18 & Чётность и делимость & Признаки делимости на 2, 3, 5 & Проверка делимости в коде\\
19 & Задачи на логику & Элементарные логические задачи & Блок-схема логической задачи\\
20 & Олимпиадная задача 1 & Задачи начального уровня & Решение в тетради\\
21 & Координатная плоскость (наглядно) & Точка и координаты & Движение спрайта по координатам\\
22 & Таймеры и темп & Понятие скорости & Переменная «время»\\
23 & Случайность & Случайное число & Блок «выдать случайное»\\
24 & Игра с угадыванием & Стратегия выигрыша & Логика игры «угадай»\\
25 & Комбинаторика & Правило произведения, начало & Генерация комбинаций\\
26 & Задачи на комбинаторику & Сколько вариантов & Визуализация комбинаций\\
27 $\star$ & Мини-проект 2 & --- & Генератор комбинаций\\
28 & Формулировка задачи & Что такое «дано, найти, условие» & Составление алгоритма\\
29 & Отладка & Поиск ошибки & «Почини программу»\\
30 & Проверка на примерах & Верификация решения & Тест правила на разных данных\\
31 & Списки и массивы & Среднее, сумма элементов & Обработка списка данных\\
32 & Мини-олимпиада №1 & 3--4 задачи & ---\\
33 & Разбор олимпиады & Разбор решений & Обсуждение\\
34 & Проектная работа & Постановка задачи проекта & Планирование\\
35 & Защита проекта & --- & Презентация игры\\
36 $\star$ & Рефлексия & Итоги полугодия & Обратная связь от учеников\\
\midrule

\multicolumn{4}{l}{\textbf{Блок 2. Данные и логика (занятия 37--72)}}\\
\midrule
37 & Сбор данных & Что такое данные & Опрос класса в таблице\\
38 & Таблицы и структура & Строка, столбец & Таблица в Scratch\\
39 & Графики & Виды графиков & Построение графика вручную\\
40 & Столбчатая диаграмма & Сравнение величин & Диаграмма в Scratch\\
41 & Логика и И/ИЛИ & Логические связи & Составление логических схем\\
42 & Таблица истинности & Все комбинации условий & Проверка в коде\\
43 & Задачи на логику & Сложные логические задачи & Блок-схема решения\\
44 & Предикаты & Утверждение о числе & Функция-проверка\\
45 & Система координат & Точка, координаты & Лабиринт по координатам\\
46 & Отрицательные числа & Знак числа & Движение «в negatives»\\
47 & Сложение отрицательных & Правила знаков & Модель правила в коде\\
48 & Дроби наглядно & Доля, часть, целое & Разрезание пиццы в коде\\
49 & Дроби и модели & Модель дроби & Сравнение долей\\
50 & Проценты наглядно & Процент как сотая часть & Скидки в игре\\
51 & Проценты & Задачи на проценты & Калькулятор скидки\\
52 & Симметрия & Симметричные фигуры & Симметрия в графике\\
53 & Геометрические фигуры & Площади фигур & Рисование фигур\\
54 & Периметр и площадь & Формулы, применение & Подсчёт площади в коде\\
55 & Клетчатая бумага & Площадь по клеткам & Построение на сетке\\
56 & Задачи на комбинаторику & Правило сложения & Выбор вариантов\\
57 & Графы наглядно & Вершины и рёбра & Построение графа\\
58 & Пути в графе & Кратчайший путь & Поиск пути\\
59 & Булевы выражения & Логика высказываний & Сложные условия в коде\\
60 $\star$ & Мини-проект 3 & --- & Игра с логическими условиями\\
61 & Обобщение данных & Как свести данные к выводу & Вывод в игре\\
62 & Столбчатая диаграмма & Чтение данных & Собственная диаграмма\\
63 & Круговая диаграмма & Доли целого & Круговая диаграмма\\
64 & Олимпиадная задача 2 & Задачи уровня района & ---\\
65 & Разбор задачи 2 & Разбор решений & Обсуждение\\
66 & Объекты и свойства & Признаки объектов & Объекты с признаками в Scratch\\
67 & Сортировка & Как сортировать & Алгоритм сортировки в игре\\
68 & Классификация & Что такое класс & Игра «Раздели на группы»\\
69 & Простые модели & Модель как правило & Правило «больше---меньше»\\
70 & Проверка модели & Как проверить правило & Тест модели на примерах\\
71 & Обобщение на новых данных & Перенос модели & ---\\
72 $\star$ & Рефлексия блока & Итоги & Обратная связь\\
\midrule

\multicolumn{4}{l}{\textbf{Блок 3. Проекты и соревнования (занятия 73--108)}}\\
\midrule
73 & Планирование проекта & Постановка задачи &  Выбор темы, распределение\\
74 & Декомпозиция & Разбиение на части &  План действий\\
75 & Сборка: часть 1 & &  Программирование\\
76 & Сборка: часть 2 & &  Программирование\\
77 & Сборка: часть 3 & &  Программирование\\
78 & Тестирование & &  Проверка на разных данных\\
79 & Отладка проекта & &  Исправление ошибок\\
80 & Улучшение & Сравнение версий &  Доработка\\
81 & Оформление & & Скриншоты, описание\\
82 & Подготовка защиты & Доказательство результата & План защиты\\
83 & Защита проекта №1 & &  Презентация (3--5 мин)\\
84 & Вопросы и обсуждение & &  Разбор вопросов\\
85 & Олимпиадная задача 3 & Сложные задачи & ---\\
86 & Разбор задачи 3 & Разбор &  ---\\
87 & Комбинаторика (продвинутая) & Правило произведения &  Подсчёт вариантов\\
88 & Принцип Дирихле & Олимпиадный приём &  Решение задач\\
89 & Игры и стратегии & Выигрышные стратегии &  Игра «выбери ход»\\
90 & Олимпиадная задача 4 & Сложные задачи & ---\\
91 & Разбор задачи 4 & Разбор &  ---\\
92 & Логика высокого уровня & Сложные цепочки условий &  Построение в коде\\
93 & Математическое моделирование & Модель реальной задачи &  Построение модели\\
94 & Модель и прогноз & Предсказание &  Прогноз в игре\\
95 & Машина как исполнитель алгоритма & Понятие автомата &  Игра «компьютер»\\
96 & Символьные вычисления & Работа с буквами &  Шифр Цезаря\\
97 & Криптография наглядно & Шифрование &  Своя шифровка\\
98 & Мини-олимпиада №2 & 4--5 задач & ---\\
99 & Разбор мини-олимпиады & Разбор &  ---\\
100 & Повторение: логика & Повторение года &  ---\\
101 & Повторение: арифметика & Повторение года &  ---\\
102 & Повторение: геометрия & Повторение года &  ---\\
103 & Экзамен-тренажёр & Пробная работа &  ---\\
104 & Проектная неделя: план & &  Разработка\\
105 & Проектная неделя: работа & &  Разработка\\
106 & Проектная неделя: работа & &  Разладка, тест\\
107 & Защита проекта №2 & &  Презентация (5--7 мин)\\
108 & Итоги года & График достижений &  Награждение, план\\
\bottomrule
\end{longtable}

\subsection{Год 2. Тематический план}

\begin{longtable}{C{1.0cm}L{4.4cm}L{4.0cm}L{4.4cm}}
\toprule
\textbf{№} & \textbf{Тема} & \textbf{Математика} & \textbf{Практика (Scratch/ML)}\\
\midrule
\endfirsthead

\toprule
\textbf{№} & \textbf{Тема} & \textbf{Математика} & \textbf{Практика (Scratch/ML)}\\
\midrule
\endhead

\multicolumn{4}{l}{\textbf{Блок 1. Алгоритмы и структуры (занятия 1--36)}}\\
\midrule
1--5 & Переменные: углубление & Типы величин, приведение & Глобальные и локальные переменные\\
6--10 & Списки: операции & Множества, элементы & Сортировка, вставка, удаление\\
11--15 & Процедуры & Повторяющиеся действия & Выделение повторов в процедуру\\
16--20 & Функции и возвраты & Функция как зависимость & Функции с параметрами\\
21--25 & Вложенные структуры & Сложные условия & Комбинации списков и функций\\
26--30 & Обработка текста & Символы, строки & Работа со строками, парсинг\\
31--34 & Рекурсия (простое) & Понятие рекурсии & Факториал рекурсивно\\
35 $\star$ & Повторение блока & --- & Игра-тест\\
36 & Проект-разминка & --- & Мини-игра «Машина времени»\\
\midrule

\multicolumn{4}{l}{\textbf{Блок 2. Данные и модели (занятия 37--72)}}\\
\midrule
37--41 & Сбор и очистка данных & Данные, шум & Фильтрация списка\\
42--46 & Статистика: среднее, мода & Среднее, медиана, мода & Считаем в коде\\
47--51 & Разброс и разброс данных & Размах, дисперсия (наглядно) & График распределения\\
52--56 & Правила и закономерности & Поиск закономерностей & Поиск правила в данных\\
57--61 & Простые модели «если--то» & Модель как правило & Классификатор «больше---меньше»\\
62--66 & Дерево решений & Разбиение на группы & Построение дерева\\
67--70 & Точность модели & Как измерить правильность & Счёт правильных ответов\\
71 $\star$ & Рефлексия блока & --- & Обсуждение\\
72 & Проект-полугодовой & --- & Защита\\
\midrule

\multicolumn{4}{l}{\textbf{Блок 3. Творчество и олимпиады (занятия 73--108)}}\\
\midrule
73--78 & Математическая графика & Построения, узоры & Фракталы простые (паутина, дерево)\\
79--84 & Криптография и коды & Шифры, кодирование & Своя шифровка, «Код Брюч»\\
85--90 & Олимпиадные приёмы & Принцип Дирихле, крайний & 5--6 задач на приём\\
91--96 & Подготовка к олимпиадам & Задачи уровня региона & Тренировочные туры\\
97--100 & Олимпиадные туры & --- & 3 пробных тура\\
101--104 & Проектная работа & Постановка & Разработка\\
105--107 & Защита проекта & --- & Презентация, вопросы\\
108 & Итоги года & График & Награждение\\
\bottomrule
\end{longtable}

\subsection{Год 3. Тематический план}

\begin{longtable}{C{1.0cm}L{4.4cm}L{4.0cm}L{4.4cm}}
\toprule
\textbf{№} & \textbf{Тема} & \textbf{Математика} & \textbf{Практика (Scratch/ML)}\\
\midrule
\endfirsthead

\toprule
\textbf{№} & \textbf{Тема} & \textbf{Математика} & \textbf{Практика (Scratch/ML)}\\
\midrule
\endhead

\multicolumn{4}{l}{\textbf{Блок 1. Систематизация (занятия 1--36)}}\\
\midrule
1--6 & Система счисления & Двоичная система, перевод & Двоичная логика в коде\\
7--12 & Отрицательные и обратные & Обратная величина & Инверсия в Scratch\\
13--18 & Формулы и функции & Формула как правило & Функция по формуле\\
19--24 & Задачи на проценты & Проценты, сложный процент & Калькулятор на проценты\\
25--30 & Логические задачи & Логика высказываний & Сложные условия\\
31--35 & Алгоритмы сортировки & Упорядочивание & Пузырьковая сортировка в Scratch\\
36 $\star$ & Повторение & --- & Тест\\
\midrule

\multicolumn{4}{l}{\textbf{Блок 2. Проекты и исследование (занятия 37--72)}}\\
\midrule
37--42 & Принцип Дирихле & Обобщение на 6--8 класс & Задачи\\
43--48 & Комбинаторика & Правила, графы & Подсчёт путей\\
49--54 & Статистика вероятностей & Вероятность наглядно & Модель «как угадать»\\
55--60 & Тригонометрия наглядно & Угол, синус (качественно) & Движение по волне\\
61--66 & Кодирование и модели & Модель как описание & Строим «модель явления»\\
67--71 & Собственное исследование & --- & Мини-исследование\\
72 $\star$ & Защита & --- & Презентация\\
\midrule

\multicolumn{4}{l}{\textbf{Блок 3. Олимпиады и итог (занятия 73--108)}}\\
\midrule
73--78 & Олимпиадные приёмы & Экстремальный, разбор случаев & Задачи\\
79--84 & Олимпиадные приёмы 2 & Инварианты, раскраска & Задачи\\
85--90 & Тренировочные туры & --- & 3 тура\\
91--96 & Пробная олимпиада & --- & Работа\\
97--100 & Разбор & --- & Разбор\\
101--105 & Финальный проект & Постановка, реализация & Своя игра или исследование\\
106--107 & Защита проекта & --- & Презентация 7--10 мин\\
108 & Итоги года и выпуск & Итоговая рефлексия & Награждение\\
\bottomrule
\end{longtable}

% ====================================================================
\newpage
\section{Траектория Б. Средние классы (5--7), 1{,}5 часа}

\subsection{Общие положения траектории}

\begin{itemize}
  \item \textbf{Возраст:} 11--14 лет. Занятие --- 90 минут.
  \item \textbf{Главная задача:} переход от блочного программирования к
  текстовому коду (Python) и одновременно --- серьёзная алгебра и
  комбинаторика.
  \item \textbf{Среда:} Python в браузере (Google Colab или Replit) ---
  не требует установки, работает на школьном компьютере.
  \item \textbf{Математика:} 5 класс --- дроби и проценты;
  6 класс --- отношения, пропорции, отрицательные числа;
  7 класс --- начала алгебры, линейные уравнения, комбинаторика.
  \item \textbf{Методический риск:} упирание в синтаксис Python.
  Решается жёстким правилом: 20\% времени на синтаксис, 80\% --- на смысл.
  \item \textbf{Мотивация:} в 6--7 классах --- подростковый бунт против
  «детского» обучения. Это шанс: дать настоящий, «взрослый» инструмент.
\end{itemize}

\subsection{Год 1. Тематический план}

\begin{longtable}{C{0.9cm}L{4.3cm}L{4.1cm}L{4.5cm}}
\toprule
\textbf{№} & \textbf{Тема} & \textbf{Математика} & \textbf{Практика (Python)}\\
\midrule
\endfirsthead

\toprule
\textbf{№} & \textbf{Тема} & \textbf{Математика} & \textbf{Практика (Python)}\\
\midrule
\endhead

\multicolumn{4}{l}{\textbf{Блок 1. Дроби, проценты и первые переменные (занятия 1--36)}}\\
\midrule
1 & Переменные и типы & Целые, дробные & \ml{print(), переменные, типы}\\
2 & Ввод и вывод & Запись результата & \ml{input(), форматирование}\\
3 & Арифметика & Порядок действий, скобки & \ml{+, -, *, /, //, \%}\\
4 & Дроби: запись & Обыкновенные дроби & \ml{Дробь через числитель/знаменатель}\\
5 & НОД и сокращение & Алгоритм Евклида & \ml{Функция НОД}\\
6 & НОК и приведение & Общий знаменатель & \ml{Функция НОК}\\
7 & Арифметика дробей & Сложение дробей & \ml{Сложение через общий знаменатель}\\
8 & Умножение и деление дробей & Действия с дробями & \ml{Умножение дробей}\\
9 & Десятичные дроби & Перевод в десятичную & \ml{float, округление}\\
10 $\star$ & Повторение & Дроби целиком & Тест-разбор\\
11 & Проценты: определение & Процент как сотая часть & \ml{Вычисление процентов}\\
12 & Проценты в задачах & Увеличение на N\% & \ml{Увеличить число на процент}\\
13 & Скидки и наценки & Скидка, наценка & \ml{Калькулятор скидок}\\
14 & Сложный процент & Рост и снижение & \ml{Рост вклада}\\
15 & Процентные отношения & Две величины в процентах & \ml{Пропорция}\\
16 & Условия & Логические условия & \ml{if / elif / else}\\
17 & Логические операции & И, ИЛИ, НЕ & \ml{and, or, not}\\
18 & Вложенные условия & Сложные ветвления & \ml{Калькулятор скидок по условию}\\
19 & Цикл for & Счёт, сумма & \ml{for, range()}\\
20 & Сумма и произведение & Прогрессия & \ml{Сумма чисел циклом}\\
21 & Цикл while & & Повторение до условия  \ml{while}\\
22 & Вложенные циклы & Таблица умножения & \ml{Таблица умножения}\\
23 & Списки & Список, элемент, индекс & \ml{Создание списка, доступ}\\
24 & Методы списка & --- & \ml{append, len, sum, max, min}\\
25 & Простые задачи на списки & --- & \ml{Обработка списка чисел}\\
26 & Функции & Определение функции & \ml{Ключевое слово def}\\
27 & Параметры и возврат & Функция как зависимость & \ml{Параметры, return}\\
28 & Области видимости & --- & \ml{Глобальные и локальные}\\
29 & Модули & --- & \ml{import math, random}\\
30 & Число пи и округления & $\pi$, округление & \ml{math.pi}\\
31 & Работа с файлами & Запись данных & \ml{Запись и чтение файла}\\
32 $\star$ & Мини-проект 1 & --- & Много «калькуляторов»\\
33 & Повторение блока & --- & Тест-разбор\\
34 & Итоговое занятие 1 & --- & Олимпиада-разминка\\
35 & --- & --- & ---\\
36 $\star$ & Рефлексия & Итоги & Обратная связь\\
\midrule

\multicolumn{4}{l}{\textbf{Блок 2. Пропорции, отношения и модели (занятия 37--72)}}\\
\midrule
37 & Отношения & Отношение, часть, целое & \ml{Функция ``часть от целого''}\\
38 & Пропорции & Основное свойство & \ml{Проверка пропорции}\\
39 & Задачи на пропорции & --- & \ml{Скупание---продажа}\\
40 & Масштаб & --- & \ml{Перевод масштаба}\\
41 & Прямо пропорциональность & Линейная зависимость & \ml{Прямая $y = kx$}\\
42 & Обратная пропорция & --- & \ml{Обратная зависимость}\\
43 & Построение графиков & График функции & \ml{Графики без библиотек}\\
44 & Линейные уравнения & Решение уравнения & \ml{Алгоритм решения $ax=b$}\\
45 & Построение графика уравнения & --- & \ml{Программа рисует прямую}\\
46 & Положение декартовой системы & Координаты, знаки & \ml{Работа с координатами}\\
47 & Расстояние между точками & Формула расстояния & \ml{Функция расстояния}\\
48 & Сложные задачи на движение & Текстовые задачи & \ml{Модель «движение»}\\
49 & Отрицательные числа & Действия с отрицательными & \ml{Числа со знаком}\\
50 & Модуль числа & --- & \ml{abs()}\\
51 & Координатная плоскость & Движения по координатам & \ml{Анимация движения}\\
52 & Параллельные и перпендикулярные & --- & \ml{Рисование фигур}\\
53 & Графики на координатной плоскости & --- & \ml{Рисуем график в коде}\\
54 & Задачи на движение (скорость) & Скорость, путь, время & \ml{Модель скорости}\\
55 & Проценты в движении & --- & \ml{Скидка + движение}\\
56 & Олимпиадная задача 1 & Задача уровня школы & ---\\
57 & Разбор задачи 1 & --- & ---\\
58 & Комбинаторика: правило сложения & --- & \ml{Генерация комбинаций}\\
59 & Комбинаторика: правило умножения & --- & \ml{Подсчёт путей}\\
60 & Размещения и перестановки & --- & \ml{Перестановки}\\
61 & Взвешенный подсчёт & --- & \ml{Взвешивание вариантов}\\
62 & Начало статистики & Сбор данных, таблица & \ml{Массив данных}\\
63 & Среднее арифметическое & Формула среднего & \ml{Вычисление среднего}\\
64 & Мода и медиана & --- & \ml{Мода, медиана}\\
65 & Размах & --- & \ml{Размах}\\
66 & Случайные числа & Случайный выбор & \ml{random}\\
67 & Моделирование: «Подбрасывание монеты» & Вероятность интуитивно & \ml{Симуляция 1000 бросков}\\
68 & Олимпиадная задача 2 & Задача уровня района & ---\\
69 & Разбор задачи 2 & --- & ---\\
70 & Мини-проект 2 & --- & Своя мини-игра или расчёт\\
71 $\star$ & Повторение блока & --- & Тест-разбор\\
72 & Рефлексия & Итоги & Обратная связь\\
\midrule

\multicolumn{4}{l}{\textbf{Блок 3. Алгебра и первая модель (занятия 73--108)}}\\
\midrule
73 & Способы задания последовательности & Ряд, арифметическая прогрессия & \ml{Закономерности}\\
74 & Арифметическая прогрессия & Формула $a_n$ & \ml{Прогрессия в коде}\\
75 & Сумма первых N & Формула суммы & \ml{Сумма прогрессии}\\
76 & Признак делимости на 3, 9 & Сумма цифр & \ml{Проверка делимости}\\
77 & Линейная функция $y = kx+b$ & Угловой коэффициент & \ml{Построение $y=kx+b$}\\
78 & Квадратные уравнения (предпросмотр) & Дискримиант & \ml{Программа-решатель}\\
79 & Координатный метод & Прямая и система & \ml{Построение двух прямых}\\
80 & Пересечение прямых & Точка пересечения & \ml{Поиск точки пересечения}\\
81 & Прямоугольный треугольник & Теорема Пифагора & \ml{Расстояние между точками}\\
82 & Проценты и банк & Вклады, проценты & \ml{Финансовая модель}\\
83 & Представление данных в коде & --- & \ml{Список, pandas позже}\\
84 & Первая «модель»: угадай число & Биномиальная структура & \ml{Алгоритм бинарного поиска}\\
85 & Что такое модель (формально) & Функция как модель & \ml{Функция-предиктор}\\
86 & Обучение на примерах (интуиция) & --- & \ml{Подбор коэффициента вручную}\\
87 & Оценка качества модели & Средняя ошибка & \ml{Подсчёт MAE}\\
88 & Первая визуализация & График точек & \ml{Рисуем данные и модель}\\
89 & Улучшение модели & Как сделать точнее & \ml{Подбор параметров}\\
90 & Олимпиадная задача 3 & --- & ---\\
91 & Разбор задачи 3 & --- & ---\\
92 & Модуль, деление на остаток & Остаток от деления & \ml{Оператор \%}\\
93 & Многоугольники & Площади фигур & \ml{Расчёт площади}\\
94 & Окружность & Формулы & \ml{Построение окружности}\\
95 & Объёмы & Формулы объёма & \ml{Модель объёма}\\
96 & Моделирование «Солнце-Земля» & --- & \ml{Анимация движения}\\
97 & Мини-проект 3 & --- & ---\\
98 & --- & --- & ---\\
99 & --- & --- & ---\\
100 & Олимпиадная задача 4 & --- & ---\\
101 & Разбор задачи 4 & --- & ---\\
102--104 & Проект «Дача данных» & Анализ данных & \ml{Собрать, обработать, визуализировать}\\
105--107 & Защита проекта & --- & Презентация, вопросы\\
108 & Итоги года & График & Награждение, план\\
\bottomrule
\end{longtable}

\subsection{Год 2. Тематический план}

\begin{longtable}{C{0.9cm}L{4.3cm}L{4.1cm}L{4.5cm}}
\toprule
\textbf{№} & \textbf{Тема} & \textbf{Математика} & \textbf{Практика (Python)}\\
\midrule
\endfirsthead

\toprule
\textbf{№} & \textbf{Тема} & \textbf{Математика} & \textbf{Практика (Python)}\\
\midrule
\endhead

\multicolumn{4}{l}{\textbf{Блок 1. Алгебра и структуры данных (занятия 1--36)}}\\
\midrule
1--4 & Раскрытие скобок & Алгебраические преобразования & \ml{Алгоритмическая проверка}\\
5--8 & Линейные уравнения & Решение, параметры & \ml{Универсальный решатель}\\
9--12 & Системы уравнений & Подстановка, сложение & \ml{Решение системы}\\
13--16 & Пропорции и разные единицы & --- & \ml{Конвертер единиц}\\
17--20 & Проценты в banking & Сложные проценты & \ml{Финансовый калькулятор}\\
21--24 & Списки: срезы & --- & \ml{Слайсы, фильтрация}\\
25--28 & Вложенные списки & Таблицы & \ml{Таблица как список списков}\\
29--32 & Строки и разбор & --- & \ml{split, join, форматирование}\\
33--35 & Функции высшего порядка & --- & \ml{map, filter, lambda}\\
36 $\star$ & Повторение & --- & Тест\\
\midrule

\multicolumn{4}{l}{\textbf{Блок 2. Статистика и модели (занятия 37--72)}}\\
\midrule
37--41 & Выборка и генерализация & Что такое выборка & \ml{Генерация данных}\\
42--46 & Среднее и дисперсия & Формулы & \ml{Расчёт, проверка}\\
47--51 & Линейная регрессия & Формула прямой & \ml{Линейная регрессия вручную}\\
52--56 & Качество регрессии & Ошибка, $R^2$ (качественно) & \ml{Подсчёт ошибок}\\
57--61 & Логистическая регрессия (интуиция) & Вероятность, порог & \ml{Пороговая классификация}\\
62--66 & Правило Байеса (простое) & Условная вероятность & \ml{Байес в коде}\\
67--70 & Задачи на теорию вероятностей & --- & \ml{Модель случайного события}\\
71 $\star$ & Повторение & --- & Тест-разбор\\
72 & Проект-полугодовой & --- & Защита\\
\midrule

\multicolumn{4}{l}{\textbf{Блок 3. Геометрия, алгоритмы, олимпиады (занятия 73--108)}}\\
\midrule
73--77 & Подобие & Коэффициент подобия & \ml{Масштабирование фигуры}\\
78--82 & Тригонометрия & Синус, косинус, тангенс & \ml{Расчёт углов, координаты}\\
83--87 & Площади и объёмы & Формулы & \ml{Модель вычислений}\\
88--92 & Комбинаторика (углублённая) & Графы, пути & \ml{Поиск пути (волновой алгоритм)}\\
93--96 & Олимпиадные приёмы & Инвариант, крайний & Задачи\\
97--99 & Тренировочные туры & --- & Туры\\
100--105 & Финальный проект & --- & Собственный ML-проект\\
106--107 & Защита & --- & Презентация\\
108 & Итоги & График & Награждение\\
\bottomrule
\end{longtable}

\subsection{Год 3. Тематический план}

\begin{longtable}{C{0.9cm}L{4.3cm}L{4.1cm}L{4.5cm}}
\toprule
\textbf{№} & \textbf{Тема} & \textbf{Математика} & \textbf{Практика (Python/ML)}\\
\midrule
\endfirsthead

\toprule
\textbf{№} & \textbf{Тема} & \textbf{Математика} & \textbf{Практика (Python/ML)}\\
\midrule
\endhead

\multicolumn{4}{l}{\textbf{Блок 1. Подготовка к ОГЭ + основы ML (занятия 1--36)}}\\
\midrule
1--5 & Линейные уравнения и неравенства & ОГЭ-темы & \ml{Решатель}\\
6--10 & Квадратные уравнения & Дискримиант, формулы & \ml{Решатель}\\
11--15 & Функции и графики & Свойства функций & \ml{Построение графиков}\\
16--20 & Прогрессии & Арифметическая, геометрическая & \ml{Моделирование}\\
21--25 & Планиметрия: база & Площади, углы, теоремы & \ml{Расчёты}\\
26--30 & Статистика и вероятность & Ожидание, частота, вероятность & \ml{Симуляции, статистика}\\
31--35 & Принцип работы линейной модели & Почему прямая --- модель & \ml{Строим свою линейную модель}\\
36 $\star$ & Повторение & --- & Тест\\
\midrule

\multicolumn{4}{l}{\textbf{Блок 2. Реальные ML-модели (занятия 37--72)}}\\
\midrule
37--41 & Классификация: идея & Что значит «обучить классификатор» & \ml{Первая классификация}\\
42--46 & Признаки (features) & Что подаём на вход & \ml{Инженерия признаков}\\
47--51 & Дерево решений & Разбиение по условиям & \ml{Строим дерево}\\
52--56 & Оценка точности & Точность, ошибки I и II типа & \ml{Матрица ошибок}\\
57--61 & Переобучение & Почему модель запоминает & \ml{Демонстрация переобучения}\\
62--66 & Разделение выборки & Обучающая и тестовая & \ml{train\_test\_split}\\
67--70 & Олимпиадные задачи & --- & Задачи\\
71 $\star$ & Повторение & --- & Тест\\
72 & Проект-полугодовой & --- & Защита мини-проекта\\
\midrule

\multicolumn{4}{l}{\textbf{Блок 3. Углубление и олимпиады (занятия 73--108)}}\\
\midrule
73--77 & Векторы и матрицы & Задачи & \ml{Работа с массивами}\\
78--82 & Логарифмы & Логарифмы, свойства & \ml{Логарифм в коде}\\
83--87 & Нейросети: интуиция & --- & \ml{Первая нейросеть (Keras)}\\
88--92 & Переобучение (углублённо) & --- & \ml{Регуляризация}\\
93--96 & Олимпиадные приёмы & --- & Задачи\\
97--99 & Тренировочные туры & --- & Туры\\
100--105 & Финальный ML-проект & --- & Свой проект с данными\\
106--107 & Защита & --- & Презентация, вопросы\\
108 & Итоги года & График & Награждение, план\\
\bottomrule
\end{longtable}

% ====================================================================
\newpage
\section{Траектория В. Старшие классы (8--9), 1{,}5 часа}

\subsection{Общие положения траектории}

\begin{itemize}
  \item \textbf{Возраст:} 13--15 лет. Занятие --- 90 минут.
  \item \textbf{Главная задача:} подготовить ученика к ОГЭ и показать,
  что школьная математика --- это основа настоящего ML.
  \item \textbf{Среда:} Python + Jupyter (Colab), pandas, numpy,
  matplotlib, scikit-learn, позже --- Keras.
  \item \textbf{Математика:} 8 класс --- алгебра (квадратные уравнения,
  неравенства), планиметрия, начала теории вероятностей;
  9 класс --- прогрессии, функции, полная планиметрия, статистика.
  \item \textbf{ML-фокус:} не «как запустить код», а «почему модель
  работает и когда не работает».
  \item \textbf{Методический риск:} недооценка теории ради «программирования».
  Правило: \textbf{каждые 2 занятия ML обязательно имеют 1 занятие чистой
  математики}. Это фундамент и будущая олимпиадная база.
  \item \textbf{Профиль:} подготовка к ОГЭ + олимпиадам + собственные ML-проекты.
\end{itemize}

\subsection{Год 1. Тематический план}

\begin{longtable}{C{0.9cm}L{4.3cm}L{4.1cm}L{4.5cm}}
\toprule
\textbf{№} & \textbf{Тема} & \textbf{Математика} & \textbf{Практика (Python/ML)}\\
\midrule
\endfirsthead

\toprule
\textbf{№} & \textbf{Тема} & \textbf{Математика} & \textbf{Практика (Python/ML)}\\
\midrule
\endhead

\multicolumn{4}{l}{\textbf{Блок 1. Алгебра --- фундамент (занятия 1--36)}}\\
\midrule
1 & Повторение: дроби и проценты & --- & \ml{Быстрый счёт}\\
2 & Раскрытие скобок, тождества & --- & \ml{Проверка тождества в коде}\\
3 & Квадратное уравнение & Дискримиант & \ml{Решатель}\\
4 & Теорема Виета & Связь корней и суммы & \ml{Проверка по Виету}\\
5 & Рациональные уравнения & --- & \ml{Универсальный решатель}\\
6 & Неравенства с модулем & Модуль, знаки & \ml{Решение неравенств}\\
7 & Системы уравнений & --- & \ml{Метод Гаусса / численно}\\
8 & Функция: область, график & Свойства & \ml{Построение}\\
9 & Квадратичная функция & Вершина, ветви & \ml{Парабола в коде}\\
10 $\star$ & Повторение блока 1 & --- & Тест-разбор\\
11 & Степень и показатель & Степени, корни & \ml{Степени в Python}\\
12 & Logarithm & Логарифмы, свойства & \ml{Logarithm в коде}\\
13 & Показательная функция & --- & \ml{График показательной}\\
14 & Логарифмическая функция & --- & \ml{График логарифмической}\\
15 & Прямая и обратная пропорция & $y = kx$ & \ml{Линейная зависимость}\\
16 & Смешанные зависимости & --- & \ml{Полиномиальные признаки}\\
17 & Вектор & Сумма, разность, скалярное произведение & \ml{Векторы в numpy}\\
18 & Матрица & Умножение, размеры & \ml{Матрицы в numpy}\\
19 & Задачи на векторы & --- & \ml{Геометрия через векторы}\\
20 & Координатная плоскость & --- & \ml{Визуализация точек}\\
21 & Расстояние между точками & --- & \ml{Функция расстояния}\\
22 & Расстояние от точки до прямой & --- & \ml{Реализация}\\
23 & Планиметрия: треугольники & Признаки подобия & \ml{Расчёты}\\
24 & Планиметрия: площади & Формулы & \ml{Площади в коде}\\
25 & Окружность & Вписанные углы & \ml{Построение}\\
26 & Олимпиадная задача 1 & --- & ---\\
27 & Разбор задачи 1 & --- & ---\\
28 & Повторение + практикум & --- & Смешанные задачи\\
29 & Мини-проект 1 & --- & Свой мини-проект\\
30 & Защита мини-проекта 1 & --- & Презентация\\
31--35 & Углубление и сложные задачи & Олимпиадные приёмы & Задачи\\
36 $\star$ & Рефлексия & Итоги & Обратная связь\\
\midrule

\multicolumn{4}{l}{\textbf{Блок 2. Статистика --- мост к ML (занятия 37--72)}}\\
\midrule
37 & Что такое ML? & --- & \ml{Введение, первая демонстрация}\\
38 & Задача и данные & Что такое датасет & \ml{Загрузка набора данных}\\
39 & Признаки и целевая переменная & --- & \ml{Смотрим на столбцы}\\
40 & Описательная статистика & Среднее, медиана, мода, дисперсия & \ml{describe()}\\
41 & Визуализация данных & Графики, распределения & \ml{matplotlib, seaborn}\\
42 & Гистограмма & --- & \ml{Гистограмма и распределение}\\
43 & Выбросы & --- & \ml{Поиск и обработка выбросов}\\
44 & Нормировка & Дроби, пропорции & \ml{Нормировка данных}\\
45 & Зачем нормировать & --- & \ml{Сравнение до/после}\\
46 & Линейная регрессия: идея & Уравнение $y = wx + b$ & \ml{Построение прямой}\\
47 & Обучение регрессии & Как найти $w$ и $b$ & \ml{Градиентный спуск (вручную)}\\
48 & Качество регрессии & MAE, RMSE & \ml{Подсчёт ошибок}\\
49 & Переобучение (регрессия) & --- & \ml{Демонстрация}\\
50 & Логистическая регрессия & Вероятность, логистическая функция & \ml{Классификация}\\
51 & Порог и матрица ошибок & --- & \ml{Матрица ошибок}\\
52 & Точность, полнота & Метрики & \ml{Считаем метрики}\\
53 & Дерево решений & --- & \ml{Обучение дерева}\\
54 & Случайный лес & --- & \ml{Ансамбль}\\
55 & Олимпиадная задача 2 & --- & ---\\
56 & Разбор задачи 2 & --- & ---\\
57 & Вероятность: классическое определение & --- & \ml{Подсчёт}\\
58 & Частота и вероятность & --- & \ml{Модель на Python}\\
59 & Теорема сложения & --- & \ml{Симуляция}\\
60 & Теорема умножения & --- & \ml{Симуляция}\\
61 & Формула Байеса & Условная вероятность & \ml{Байес в коде}\\
62 & Пример Байеса: тест на болезнь & --- & \ml{Спам-фильтр / тест}\\
63 & Ожидание случайной величины & Математическое ожидание & \ml{Расчёт}\\
64 & Дисперсия случайной величины & --- & \ml{Расчёт}\\
65 & Моделирование случайных явлений & --- & \ml{Модель-симуляция}\\
66 & Статистика: выборочная & --- & \ml{Выборка}\\
67 & Доверительный интервал & --- & \ml{Интервал}\\
68 & Олимпиадная задача 3 & --- & ---\\
69 & Разбор задачи 3 & --- & ---\\
70 & Мини-проект 2 & --- & ---\\
71 $\star$ & Повторение & --- & Тест\\
72 & Защита мини-проекта 2 & --- & Презентация\\
\midrule

\multicolumn{4}{l}{\textbf{Блок 3. Нейросети и финальный проект (занятия 73--108)}}\\
\midrule
73 & Нейрон: что это & --- & \ml{Модель нейрона}\\
74 & Функция активации & --- & \ml{Сигмоида, ReLU}\\
75 & Многослойная сеть & --- & \ml{Строим сеть}\\
76 & Обучение нейросети & Градиентный спуск (наглядно) & \ml{Обучение в Keras}\\
77 & Переобучение нейросети & --- & \ml{Демонстрация}\\
78 & Данные: грязные и чистые & --- & \ml{Очистка}\\
79 & Признаки: создание & --- & \ml{Feature engineering}\\
80 & Признаки: отбор & --- & \ml{Отбор признаков}\\
81 & Оценка модели: честно & --- & \ml{Cross-validation}\\
82 & Проблема утечки данных & --- & \ml{Пример утечки}\\
83 & Олимпиадная задача 4 & --- & ---\\
84 & Разбор задачи 4 & --- & ---\\
85 & Олимпиадные приёмы & --- & Задачи\\
86--88 & Тренировочные туры & --- & 3 тура\\
89--91 & Разбор туров & --- & Разбор\\
92--98 & Финальный ML-проект & --- & Свой проект\\
99--101 & Финальный ML-проект (продолжение) & --- & Отладка\\
102 & Подготовка защиты & --- & План\\
103 & Защита проекта & --- & Презентация 10 мин\\
104 & Вопросы и защита & --- & Ответы\\
105 & ОГЭ-пробник & --- & Пробная работа\\
106 & Разбор пробника & --- & Разбор\\
107 & Итоговая рефлексия & --- & ---\\
108 & Итоги года & График & Награждение\\
\bottomrule
\end{longtable}

\subsection{Год 2. Тематический план}

\begin{longtable}{C{0.9cm}L{4.3cm}L{4.1cm}L{4.5cm}}
\toprule
\textbf{№} & \textbf{Тема} & \textbf{Математика} & \textbf{Практика (Python/ML)}\\
\midrule
\endfirsthead

\toprule
\textbf{№} & \textbf{Тема} & \textbf{Математика} & \textbf{Практика (Python/ML)}\\
\midrule
\endhead

\multicolumn{4}{l}{\textbf{Блок 1. Глубокая алгебра и подготовка к ОГЭ (занятия 1--36)}}\\
\midrule
1--5 & Квадратные уравнения (углублённо) & Дискриминат, приёмы & \ml{Решатель}\\
6--10 & Неравенства и системы & --- & \ml{Решение}\\
11--15 & Степени и корни & --- & \ml{Модели}\\
16--20 & Функции и свойства & --- & \ml{Анализ графиков}\\
21--25 & Арифметическая и геометрическая прогрессии & Формулы & \ml{Задачи}\\
26--30 & Планиметрия (полная) & Площади, подобие, окружность & \ml{Расчёты}\\
31--34 & Стереометрия: основы & Объёмы, векторы & \ml{3D-модель}\\
35--36 $\star$ & Повторение & --- & Тест\\
\midrule

\multicolumn{4}{l}{\textbf{Блок 2. ML-инструменты (занятия 37--72)}}\\
\midrule
37--41 & Pandas: основы & --- & \ml{Чтение, фильтрация, группировка}\\
42--46 & Pandas: агрегации & Среднее, медиана & \ml{groupby}\\
47--51 & Matplotlib: графики & --- & \ml{Типы графиков}\\
52--56 & Очистка данных & --- & \ml{Пропуски, дубликаты}\\
57--61 & Кодирование категорий & --- & \ml{Label encoding, one-hot}\\
62--66 & Признаки: масштабирование & --- & \ml{Стандартизация}\\
67--70 & Признаки: отбор & --- & \ml{Корреляция, важность}\\
71 $\star$ & Повторение & --- & Тест\\
72 & Проект-полугодовой & --- & Защита\\
\midrule

\multicolumn{4}{l}{\textbf{Блок 3. Продвинутый ML и олимпиады (занятия 73--108)}}\\
\midrule
73--77 & Модели: SVM & --- & \ml{Линейная и ядерная}\\
78--82 & Ансамбли & --- & \ml{Бустинг}\\
83--87 & Временные ряды & --- & \ml{Прогноз}\\
88--92 & Рекомендательные системы & --- & \ml{Коллаборативная фильтрация}\\
93--96 & Олимпиадные приёмы & --- & Задачи\\
97--99 & Туры & --- & Олимпиады\\
100--105 & Финальный проект & --- & Свой ML-проект\\
106--107 & Защита & --- & Презентация\\
108 & Итоги & График & Награждение\\
\bottomrule
\end{longtable}

\subsection{Год 3. Тематический план}

\begin{longtable}{C{0.9cm}L{4.3cm}L{4.1cm}L{4.5cm}}
\toprule
\textbf{№} & \textbf{Тема} & \textbf{Математика} & \textbf{Практика (Python/ML)}\\
\midrule
\endfirsthead

\toprule
\textbf{№} & \textbf{Тема} & \textbf{Математика} & \textbf{Практика (Python/ML)}\\
\midrule
\endhead

\multicolumn{4}{l}{\textbf{Блок 1. Олимпиадная математика (занятия 1--36)}}\\
\midrule
1--4 & Принцип Дирихле & Обобщение & Задачи\\
5--8 & Экстремальный принцип & --- & Задачи\\
9--12 & Инварианты & --- & Задачи\\
13--16 & Раскраска & --- & Задачи\\
17--20 & Метод математической индукции & --- & Задачи\\
21--24 & Оценка + пример & --- & Задачи\\
25--28 & Задачи на делимость & --- & Задачи\\
29--32 & Неравенства (олимпиадные) & --- & Задачи\\
33--35 & Текстовые задачи & --- & Задачи\\
36 $\star$ & Повторение & --- & Тест\\
\midrule

\multicolumn{4}{l}{\textbf{Блок 2. Глубокий ML и исследования (занятия 37--72)}}\\
\midrule
37--41 & Глубокое обучение: CNN & --- & \ml{Свёрточные сети}\\
42--46 & Глубокое обучение: NLP & --- & \ml{Текстовые модели}\\
47--51 & Обучение без учителя & --- & \ml{Кластеризация}\\
52--56 & Снижение размерности & --- & \ml{PCA}\\
57--61 & Объяснимость (интерпретация) & --- & \ml{Важность признаков}\\
62--66 & Этика ML & --- & \ml{Смещения, приватность}\\
67--70 & Собственное исследование & --- & Мини-исследование\\
71 $\star$ & Повторение & --- & Тест\\
72 & Защита исследования & --- & Презентация\\
\midrule

\multicolumn{4}{l}{\textbf{Блок 3. Финальный проект и выпуск (занятия 73--108)}}\\
\midrule
73--78 & Олимпиадные туры & --- & Туры\\
79--82 & Разбор & --- & Разбор\\
83--90 & Финальный капитальный проект & --- & Большой ML-проект\\
91--94 & Доработка проекта & --- & Отладка\\
95 & Подготовка защиты & --- & Презентация\\
96--97 & Защита проекта & --- & Защита, вопросы\\
98--100 & ОГЭ / вступительные & --- & Пробная работа\\
101--102 & Разбор & --- & Разбор\\
103--105 & Индивидуальный разбор & --- & Работа с учеником\\
106--107 & Куда дальше & --- & Консультация, план\\
108 & Выпускной & Итоги, график & Награждение\\
\bottomrule
\end{longtable}

% ====================================================================
\newpage
\section{Система оценивания и аттестации}

\subsection{Три компонента оценки}

\begin{table}[H]
\centering
\caption{Структура итоговой оценки за год}
\label{tab:grading}
\begin{tabular}{L{4.2cm}C{2.4cm}L{7.4cm}}
\toprule
\textbf{Компонент} & \textbf{Вес} & \textbf{Что оценивается}\\
\midrule
Текущие работы & 40\% & Домашние задания, практикумы, ответы
у доски, участие\\
Олимпиады и контрольные & 30\% & Умение решать нестандартные задачи,
работа в условиях ограниченного времени\\
Проект & 30\% & Постановка задачи, реализация, защита,
ответы на вопросы\\
\bottomrule
\end{tabular}
\end{table}

\subsection{Шкала уровней}

\begin{longtable}{L{2.6cm}L{3.4cm}L{8.4cm}}
\toprule
\textbf{Уровень} & \textbf{Оценка} & \textbf{Описание}\\
\midrule
\endfirsthead

\toprule
\textbf{Уровень} & \textbf{Оценка} & \textbf{Описание}\\
\midrule
\endhead

Начальный & 2--3 & Решает типовые задачи с помощью учителя,
теряется в заданиях с неполным условием\\
Базовый & 4 & Решает типовые задачи самостоятельно,
владеет базовым синтаксисом\\
Повышенный & 5 & Решает задачи повышенной сложности,
применяет математику в нестандартных ситуациях\\
Высокий & --- & Решает олимпиадные задачи, строит
собственные модели, защищает решения\\
\bottomrule
\end{longtable}

\note{Оценка «5» не ставится за прочитанную теорию. Только за
решённую задачу или работающий проект. Это ключевой принцип
мотивации: ученик ценит не «знание», а «умение».}

\subsection{Формат обратной связи}

Каждое занятие заканчивается письменной обратной связью ---
одним предложением по каждому ученику. Формат:

\begin{quote}
\small
\textbf{Дата. Имя ученика.}
Сильная сторона: [что именно получилось].
Зона роста: [одна конкретная вещь].
Задание к следующему занятию: [конкретно].
\end{quote}

\noindent
Это занимает у преподавателя 10--15 минут после занятия, но
удерживает ученика и даёт родителям ощущение реальной работы.
Заполненные листы отдаются ученику.

\subsection{Аттестационные события}

\begin{longtable}{L{2.4cm}L{4.4cm}L{7.2cm}}
\toprule
\textbf{Когда} & \textbf{Событие} & \textbf{Что готовить}\\
\midrule
\endfirsthead

\toprule
\textbf{Когда} & \textbf{Событие} & \textbf{Что готовить}\\
\midrule
\endhead

Занятие 32 (А) & Мини-олимпиада №1 & 3--4 задачи по мотивации
без повторения\\
Занятие 56 (Б) & Олимпиадный тур & 5--6 задач уровня школы/района\\
Занятие 68 (В) & Олимпиадный тур & 5--6 задач уровня района
либо школьного этапа ВсОШ\\
Занятие 84--108 & Защита проекта & Готовый проект + защита 10 мин\\
Конец года & Итоговая работа & Смешанные задачи всех блоков\\
\end{longtable}

% ====================================================================
\section{Работа с одарёнными и отстающими}

\subsection{Правило трёх уровней сложности}

Каждое домашнее задание содержит три обязательных уровня.
Ученик выбирает свой, но рекомендуемый уровень нельзя пропускать
два раза подряд --- это правило не даёт «отсиживаться» на лёгком.

\begin{table}[H]
\centering
\caption{Три уровня заданий}
\label{tab:levels}
\begin{tabular}{L{2.6cm}L{4.2cm}L{7.4cm}}
\toprule
\textbf{Уровень} & \textbf{Тип задания} & \textbf{Для кого}\\
\midrule
Базовый & 2--3 однотипные задачи & Все, обязательно\\
Повышенный & Задача с неполным условием, требующая
размышления & Те, кто справился с базовым\\
Олимпиадный & 1 задача уровня соревнований & Сильные ученики
и все, кто хочет расти\\
\bottomrule
\end{tabular}
\end{table}

\subsection{Отстающие ученики}

\begin{enumerate}
  \item \textbf{Правило «обязательного минимума».} Каждый ученик
  получает понятный минимум, который реально выполнить. Требование
  --- не «сделай все», а «сделай базовый уровень плюс одно задание
  повышенного за два занятия».
  \item \textbf{Дробление материала.} Большая задача делится на
  2--3 подзадачи с отдельным счётом.
  \item \textbf{Повторение вместо нового.} Если ученик не понял тему,
  следующее занятие начинается с 10-минутного повторения вместо
  нового материала.
  \item \textbf{Определённый успех.} Каждый раз отмечаем конкретный
  прогресс: «ты теперь умеешь то, что не умел две недели назад».
  \item \textbf{Никаких сравнений в классе.} Не выносить оценки
  на доску, не сравнивать учеников друг с другом публично.
\end{enumerate}

\subsection{Одарённые ученики}

\begin{enumerate}
  \item \textbf{Олимпиадный уровень с первого занятия.} Для сильных
  учеников готовить задачи уровня муниципального этапа ВсОШ
  уже в младших классах.
  \item \textbf{Роль наставника.} Сильный ученик может вести
  мини-объяснение для отстающих --- это консолидирует его знания.
  \item \textbf{Углублённые проекты.} Свой проект, исследование,
  построение нетривиальной модели.
  \item \textbf{Выступления.} Подготовка к олимпиадам, научно-популярные
  выступления, помощь младшим группам.
  \item \textbf{Персональный темп.} Одарённому ученику разрешается
  двигаться быстрее, но на каждом занятии он должен сделать
  хотя бы одно задание базового уровня, чтобы не терять связь с программой.
\end{enumerate}

% ====================================================================
\section{Методика, снижающая отток учеников}

\subsection{Почему ученики уходят}

\begin{longtable}{L{5.2cm}L{5.4cm}L{5.4cm}}
\toprule
\textbf{Причина} & \textbf{Как проявляется} & \textbf{Что делать}\\
\midrule
\endfirsthead

\toprule
\textbf{Причина} & \textbf{Как проявляется} & \textbf{Что делать}\\
\midrule
\endhead

Скучно & Пропуски, «не хочу» & Повысить уровень сложности,
дать олимпиадные задачи\\
\addlinespace
Слишком сложно & Не делает домашние задания, тихий & Правило
трёх уровней, упрощение базы, дробление\\
\addlinespace
Непонятно, зачем & «Зачем нам это?» & Показывать связь с ML,
реальными задачами, карьерой\\
\addlinespace
Плохие отношения & Избегает занятий, ссоры & Индивидуальный подход,
смена формата, разговор с родителями\\
\addlinespace
Нет видимого результата & «Я ничего не понял» & Быстрые победы,
демонстрация прогресса, график достижений\\
\addlinespace
Перегруз & Жалобы, усталость & Снизить нагрузку, обсудить с родителями,
дать перерыв\\
\addlinespace
Не нравится формат & Уходит в другую секцию & Сменить онлайн/офлайн,
изменить группу, дать выбор\\
\bottomrule
\end{longtable}

\subsection{Профилактика оттока: 12 приёмов}

\begin{enumerate}
  \item \textbf{Первое занятие --- бесплатное.} Пробная встреча снимает
  главный барьер: «а вдруг не понравится». Конверсия в оплачивающий
  ученик с бесплатного пробника --- 60--70\%.
  \item \textbf{Видимый прогресс каждые 3 занятия.} Ведение
  личной карточки достижений ученика.
  \item \textbf{Обратная связь от ученика.} Спрашивайте
  «что тебе понравилось на прошлом занятии» --- это удерживает
  вовлечённость.
  \item \textbf{Конкурсы и мини-соревнования.} Ежемесячный
  внутренний рейтинг группы, мини-олимпиады.
  \item \textbf{Групповая динамика.} Старайтесь, чтобы в группе
  были и сильные, и начинающие --- это полезно обеим сторонам.
  \item \textbf{Творческие проекты.} Каждый блок заканчивается
  чем-то, что можно показать родителям.
  \item \textbf{Связь с реальными задачами.} «Вот эту модель
  используют в медицине / в рекомендациях / в прогнозе погоды».
  \item \textbf{Совместные достижения.} Победы группы
  (совместные проекты, командные олимпиады).
  \item \textbf{Гибкий график.} Перенос занятий, записи,
  возможность подключиться онлайн вместо офлайна.
  \item \textbf{Прозрачность для родителей.} Отчёт о прогрессе,
  рекомендации, план на год --- родитель видит ценность.
  \item \textbf{«Правило трёх занятий»:} если ученик хочет уйти,
  предложите три занятия с изменённой программой
  (больше олимпиад, больше практики, другой формат).
  \item \textbf{Финальный аккорд.} Перед каждым блоком --- мини-проект
  с презентацией, чтобы у ученика был завершённый результат.
\end{enumerate}

% ====================================================================
\section{Типовые ошибки и как их предотвращать}

\begin{longtable}{L{4.2cm}L{5.4cm}L{5.4cm}}
\toprule
\textbf{Ошибка преподавателя} & \textbf{Последствия} & \textbf{Как избежать}\\
\midrule
\endfirsthead

\toprule
\textbf{Ошибка преподавателя} & \textbf{Последствия} & \textbf{Как избежать}\\
\midrule
\endhead

Объяснять теорию дольше 20 минут & Потеря внимания, непонимание & Тайминг, чередование
теории и практики\\
\addlinespace
Задавать слишком трудные ДЗ & Отстающие сдаются & Правило трёх уровней\\
\addlinespace
Игнорировать одарённых & Сильные ученики уходят & Олимпиадные задачи
каждому занятию\\
\addlinespace
Давать только практику без теории & Нет фундамента, «работает --- не понимаю» & 20\%
теории, но обязательно\\
\addlinespace
Оценивать только результат & Ученик боится ошибки & Оценка прогресса,
обратная связь\\
\addlinespace
Смешивать уровни в одной группе & Сильные скучают, слабые отстают & Правило
трёх уровней внутри группы\\
\addlinespace
Пропускать олимпиады & Нет «потолка», нет цели & Мини-олимпиада каждые 4 недели\\
\addlinespace
Страх перед ошибкой у ученика & Избегание задач, «не хочу ошибаться» & Нормализация
ошибки, разбор на занятии\\
\addlinespace
Путать математику и ML & Две несвязанные системы & Использовать
таблицу соответствий\\
\addlinespace
Перегружать ML в младших классах & Потеря интереса, «слишком сложно» & Младшие ---
только визуальные среды и аналогии\\
\addlinespace
Не следить за нагрузкой & Выгорание учеников и преподавателя & Мониторинг,
снижение нагрузки при признаках\\
\addlinespace
Не давать обратной связи & Родители не видят ценности & Письменный отчёт
каждый блок\\
\bottomrule
\end{longtable}

% ====================================================================
\newpage
\section{Чек-листы преподавателя}

\subsection{Чек-лист подготовки к занятию}

\begin{longtable}{L{0.8cm}L{7.2cm}L{1.6cm}L{5.2cm}}
\toprule
\textbf{№} & \textbf{Пункт} & \textbf{Готово} & \textbf{Комментарий}\\
\midrule
\endfirsthead

\toprule
\textbf{№} & \textbf{Пункт} & \textbf{Готово} & \textbf{Комментарий}\\
\midrule
\endhead

1 & Продумана цель занятия (одна, конкретная) & \textbullet &\\
2 & Подготовлены материалы / открыта среда & \textbullet &\\
3 & Проверено домашнее задание & \textbullet &\\
4 & Подготовлены задачи трёх уровней & \textbullet &\\
5 & Определено, какая тема математики свяжется с ML & \textbullet &\\
6 & Запланирован мини-проект или задача-вызов & \textbullet &\\
7 & Подготовлены карточки обратной связи & \textbullet &\\
8 & Запланирован тайминг (теория $\leq$ 24 мин) & \textbullet &\\
\bottomrule
\end{longtable}

\subsection{Чек-лист после занятия}

\begin{longtable}{L{0.8cm}L{7.2cm}L{1.6cm}L{5.2cm}}
\toprule
\textbf{№} & \textbf{Пункт} & \textbf{Сделано} & \textbf{Комментарий}\\
\midrule
\endfirsthead

\toprule
\textbf{№} & \textbf{Пункт} & \textbf{Сделано} & \textbf{Комментарий}\\
\midrule
\endhead

1 & Записано фактическое время по этапам & \textbullet &\\
2 & Отмечены типичные ошибки учеников & \textbullet &\\
3 & Выдана обратная связь каждому ученику & \textbullet &\\
4 & Скорректирован план следующего занятия & \textbullet &\\
5 & Отмечены ученики группы риска (отток) & \textbullet &\\
6 & Записаны идеи и улучшения & \textbullet &\\
\bottomrule
\end{longtable}

\subsection{Чек-лист месяца (для самопроверки)}

\begin{longtable}{L{0.8cm}L{7.2cm}L{1.6cm}L{5.2cm}}
\toprule
\textbf{№} & \textbf{Пункт} & \textbullet &\\
\midrule
\endfirsthead

\toprule
\textbf{№} & \textbf{Пункт} & \textbullet &\\
\midrule
\endhead

1 & Проведено $\geq$ 10 занятий & \textbullet &\\
2 & Каждый ученик получил $\geq$ 8 карточек обратной связи & \textbullet &\\
3 & Проведена мини-олимпиада & \textbullet &\\
4 & Каждый ученик решил $\geq$ 3 задачи олимпиадного уровня & \textbullet &\\
5 & Проведён мини-проект & \textbullet &\\
6 & Отмечены прогресс и зоны роста & \textbullet &\\
7 & Нет учеников на грани оттока (или есть план) & \textbullet &\\
8 & Родители получили отчёт & \textbullet &\\
\bottomrule
\end{longtable}

% ====================================================================
\section{Приложения}

\subsection{Приложение А. Шаблон карточки обратной связи}

\begin{table}[H]
\centering
\fbox{\begin{minipage}{0.9\textwidth}
\textbf{КАРТОЧКА УЧЕНИКА} \hfill Дата: \rule{2cm}{0.4pt}\\
\hrulefill
\textbf{Имя:} \rule{5cm}{0.4pt} \hfill \textbf{Занятие:} \rule{1.5cm}{0.4pt}\\[0.3cm]
\textbf{Тема занятия:} \rule{9cm}{0.4pt}\\[0.5cm]
\textbf{Сильная сторона:} \rule{11cm}{0.4pt}\\[0.5cm]
\textbf{Зона роста:} \rule{11cm}{0.4pt}\\[0.5cm]
\textbf{Домашнее задание:} \rule{11cm}{0.4pt}\\[0.5cm]
\hrulefill
\textbf{Оценка:} $\square$3\quad$\square$4\quad$\square$5\quad
\textbf{Уровень:} $\square$база\quad$\square$повышенный
\end{minipage}}
\end{table}

\subsection{Приложение Б. Топ-20 задач-вызовов для младших классов}

\begin{enumerate}
  \item Лампочки: 10 лампочек, выключатель на каждую, как за $2$ действия
  определить, что работает?
  \item Тремя прямыми разрезать квадрат на 4 части.
  \item Поезд идёт 3 часа со скоростью 80 км/ч. Какое расстояние?
  \item Есть 8 монет, одна легче. Как найти за 2 взвешивания?
  \item Сколько треугольников в фигуре из 6 отрезков?
  \item Пять друзей пожали руки друг другу. Сколько рукопожатий?
  \item Бабушка испекла 12 пирожков, каждый разрезали на 2. Сколько кусков?
  \item Сколько двухзначных чисел, кратных 7?
  \item Расставьте 8 фишек в ряд так, чтобы между соседними было
  разное число пустых мест.
  \item Число 2026 разбить на 2026 слагаемых так, чтобы сумма
  делилась на 5.
  \item Пять расставленных часов показывают разное время. Какое
  время они показывают?
  \item Из 9 точек соединить все, чтобы было ровно 20 отрезков.
  \item Мальчик встал в очередь. Перед ним 3 человека, за ним
  ещё 5. Сколько всего?
  \item Три прямые пересекаются. На сколько частей делится плоскость?
  \item Сколько существует трёхзначных чисел, у которых цифры
  идут в возрастающем порядке?
  \item Ваня забыл расписание. Расставьте звонки 1--7 в случайном
  порядке. Какова вероятность?
  \item Пять одинаковых палочек, как из них составить 3 квадрата?
  \item Прямоугольник $2\times3$ разрезали на 6 единичных квадратиков.
  Сколько различных разрезаний?
  \item Если завтра в 10 утра будет дождь, через сколько часов
  будет дождь? (правильный ответ: через 2 часа)
  \item Нарисуйте фигуру из 12 точек и 8 отрезков, чтобы
  получился один многоугольник без самопересечений.
\end{enumerate}

\subsection{Приложение В. Типовые проекты}

\begin{longtable}{L{2.2cm}L{4.2cm}L{3.4cm}L{4.2cm}}
\toprule
\textbf{Класс} & \textbf{Проект} & \textbf{Математика} & \textbf{ML / программирование}\\
\midrule
\endfirsthead

\toprule
\textbf{Класс} & \textbf{Проект} & \textbf{Математика} & \textbf{ML / программирование}\\
\midrule
\endhead

2--3 & «Игра-угадайка» & Угадывание числа,
счёт & Условия и циклы в Scratch\\
\addlinespace
4 & «Генератор комбинаций» & Правило произведения & Вложенные циклы,
списки\\
\addlinespace
5 & «Калькулятор скидок» & Проценты, условия & Python: условия,
функции\\
\addlinespace
6 & «Симуляция бросков монеты» & Вероятность, среднее & Моделирование,
random\\
\addlinespace
7 & «Линейная модель» & Линейная функция, график & Построение
прямой по данным\\
\addlinespace
8 & «Предсказание цены» & Корреляция, тренд & Линейная
регрессия вручную\\
\addlinespace
8 & «Спам-фильтр» & Формула Байеса & Классификация\\
\addlinespace
9 & «Дерево решений» & Разбиение на группы & Обучение дерева\\
\addlinespace
9 & «Нейросеть» & Градиентный спуск & Keras: первая сеть\\
\addlinespace
9 & «Свой ML-проект» & Статистика & Полный цикл
данных и модели\\
\bottomrule
\end{longtable}

\subsection{Приложение Г. Список ресурсов и инструментов}

\begin{longtable}{L{3.6cm}L{5.4cm}L{5.0cm}}
\toprule
\textbf{Инструмент} & \textbf{Для чего} & \textbf{Примечание}\\
\midrule
\endfirsthead

\toprule
\textbf{Инструмент} & \textbf{Для чего} & \textbf{Примечание}\\
\midrule
\endhead

Google Colab & Python в браузере & Бесплатно, установка не нужна,
работает на школьном компьютере\\
\addlinespace
Replit & Альтернатива Colab & Онлайн-IDE, удобен для
совместной работы\\
\addlinespace
Scratch (оффлайн) & Младшие классы & Работает без интернета\\
\addlinespace
Kaggle Datasets & Наборы данных для ML & Бесплатные реальные
наборы данных\\
\addlinespace
pandas & Обработка таблиц & С 8 класса\\
\addlinespace
numpy & Вычисления & С 7 класса\\
\addlinespace
matplotlib & Графики & С 8 класса\\
\addlinespace
scikit-learn & ML-модели & С 9 класса\\
\addlinespace
Keras & Нейросети & С 9 года, позже\\
\addlinespace
Telegram & Канал, чат, бот & Основной канал связи с
учениками и родителями\\
\addlinespace
Google Формы & Тесты, анкеты & Для обратной связи\\
\addlinespace
Miro / Jamboard & Схемы, модели & Для совместной работы
над проектами\\
\bottomrule
\end{longtable}

\note{Обязательное правило: все данные для ML-проектов должны
быть либо синтетическими (сгенерированные учеником), либо
открытыми (Kaggle, gov data). Персональные данные реальных
людей без согласия использовать нельзя --- это и юридически,
и этически важно.}

\subsection{Приложение Д. Сводная таблица всего плана}

\begin{longtable}{L{4.4cm}C{3.0cm}C{3.0cm}C{3.0cm}}
\toprule
\textbf{Траектория} & \textbf{Год 1} & \textbf{Год 2} & \textbf{Год 3}\\
\midrule
\endfirsthead

\toprule
\textbf{Траектория} & \textbf{Год 1} & \textbf{Год 2} & \textbf{Год 3}\\
\midrule
\endhead

А. Младшие (1--4): занятий & 108 & 108 & 108\\
\quad из них часов (1 ч) & 108 & 108 & 108\\
Б. Средние (5--7): занятий & 108 & 108 & 108\\
\quad из них часов (1{,}5 ч) & 162 & 162 & 162\\
В. Старшие (8--9): занятий & 108 & 108 & 108\\
\quad из них часов (1{,}5 ч) & 162 & 162 & 162\\
\midrule
Всего занятий (3 траектории) & 324 & 324 & 324\\
Всего часов & 432 & 432 & 432\\
\midrule
\multicolumn{4}{r}{Итого за 3 года: 972 занятия, 1296 часов}\\
\bottomrule
\end{longtable}

\vspace{0.5cm}

\begin{center}
\small
\textbf{Документ подготовлен как методическое обеспечение
для курса «Математика и машинное обучение для 1--9 классов».}\\
Каждая траектория рассчитана на 324 занятия (108 в год, три года);
все три траектории --- 972 занятия и 1296 часов.\\
Планы Года 1 даны поурочно (все 108 занятий каждой траектории);
планы Года 2 и Года 3 сгруппированы по 5--6 занятий --- блоками
с единой темой. Перед началом года преподаватель разворачивает
каждый блок в поурочный план с указанием конкретных примеров и задач.
\end{center}

\end{document}
